\documentclass[10.5pt,a4paper]{report}
\usepackage[margin=2.0cm,top=2.2cm,bottom=2.2cm]{geometry}
\usepackage[T1]{fontenc}
\usepackage{lmodern}
\usepackage{microtype}
\usepackage{amsmath,amssymb}
\usepackage[dvipsnames,table]{xcolor}
\usepackage{booktabs,longtable,array,tabularx}
\usepackage{enumitem}
\usepackage{fancyhdr}
\usepackage{listings}
\usepackage[colorlinks=true,linkcolor=NavyBlue,urlcolor=NavyBlue,citecolor=NavyBlue]{hyperref}

\setlist{itemsep=1pt,topsep=3pt,parsep=0pt}
\setlength{\parindent}{0pt}
\setlength{\parskip}{4pt}
\renewcommand{\arraystretch}{1.18}
\setcounter{tocdepth}{2}
\setcounter{secnumdepth}{3}
\emergencystretch=2em

% ---------- column helpers ----------
\newcolumntype{L}[1]{>{\raggedright\arraybackslash}p{#1}}
\newcolumntype{C}[1]{>{\centering\arraybackslash}p{#1}}

% ---------- colours ----------
\definecolor{pickc}{HTML}{1B7F3B}
\definecolor{avoidc}{HTML}{B3261E}
\definecolor{trapc}{HTML}{C26A00}
\definecolor{exc}{HTML}{1F5FA8}
\definecolor{keyc}{HTML}{555555}
\definecolor{decc}{HTML}{6A2C91}
\definecolor{hdrc}{HTML}{E8EEF7}

% ---------- boxes (no external packages: lrbox + fcolorbox) ----------
\newsavebox{\cboxbin}
\setlength{\fboxsep}{6pt}
\setlength{\fboxrule}{0.8pt}
\newenvironment{cbox}[3]{%
  \par\vspace{5pt}\noindent
  \def\cboxfc{#1}\def\cboxbg{#2}%
  \begin{lrbox}{\cboxbin}\begin{minipage}{\dimexpr\linewidth-2\fboxsep-2\fboxrule\relax}
  \setlength{\parskip}{3pt}\textbf{\textcolor{#1}{#3}}\ \ignorespaces
}{%
  \end{minipage}\end{lrbox}\fcolorbox{\cboxfc}{\cboxbg}{\usebox{\cboxbin}}\par\vspace{5pt}}

\newenvironment{pick}[1][Choose it when]{\begin{cbox}{pickc}{pickc!6}{#1:}}{\end{cbox}}
\newenvironment{avoid}[1][Do NOT choose it when]{\begin{cbox}{avoidc}{avoidc!5}{#1:}}{\end{cbox}}
\newenvironment{trap}[1][Exam trap]{\begin{cbox}{trapc}{trapc!7}{#1:}}{\end{cbox}}
\newenvironment{ex}[1][Worked example]{\begin{cbox}{exc}{exc!5}{#1.}}{\end{cbox}}
\newenvironment{key}[1][Key facts]{\begin{cbox}{keyc}{black!4}{#1.}}{\end{cbox}}
\newenvironment{dec}[1][Decision rule]{\begin{cbox}{decc}{decc!5}{#1:}}{\end{cbox}}

% ---------- small helpers ----------
% source tag: which deck(s) to open in the exam
\newcommand{\src}[1]{\par\vspace{-2pt}\hfill{\footnotesize\itshape\textcolor{gray}{Slides: #1}}\par}
% material NOT stated in the slides (general knowledge / my derivation) -- verify before relying
\newcommand{\bs}{\textsuperscript{\textcolor{avoidc}{$\dagger$}}}
\newcommand{\cc}{\,cc}
\newcommand{\AMAT}{\mathrm{AMAT}}
\newcommand{\CPI}{\mathrm{CPI}}

\newcommand{\tblhead}[1]{\rowcolor{hdrc}#1}

% ---------- header/footer ----------
\pagestyle{fancy}
\fancyhf{}
\fancyhead[L]{\small\itshape\leftmark}
\fancyhead[R]{\small\itshape Processor Design Exam Playbook}
\renewcommand{\chaptermark}[1]{\markboth{\thechapter.\ #1}{}}
\fancyfoot[C]{\small\thepage}
\renewcommand{\headrulewidth}{0.3pt}

\title{\Huge\bfseries Processor Design Exam Playbook\\[6pt]
\Large Which architectural feature to choose, for which scenario, and why}
\author{Built from the Computer Architecture lecture decks\\(Dr.\ Noor Mahammad Sk + supplementary decks)}
\date{}

\begin{document}
\maketitle
\thispagestyle{empty}

\begin{abstract}
\noindent This is an open-book \emph{decision manual}. The exam will hand you a scenario (workload, power/cost/area limits, compatibility needs, latency numbers) and ask you to configure a processor and justify the choices. Every chapter therefore follows the same pattern: (1) a \textbf{decision map} (constraint $\rightarrow$ feature), (2) each feature with \textbf{what it fixes / what it costs / when it wins / when it loses}, (3) the \textbf{numbers from the slides} you can quote, (4) \textbf{exam traps}. Chapter~\ref{ch:scen} puts it all together with full worked scenarios and a reusable answer template.

\medskip
\noindent\textbf{Conventions.}
\emph{Slides:} tags at the end of a section tell you which deck to open (deck names are listed in Appendix~\ref{app:src}).
A red dagger\bs{} marks something that is \emph{not} stated in the slides (my general knowledge or my own arithmetic on slide numbers); it is correct as far as I know, but treat it as a supplement and prefer slide wording in your answer.

\medskip
\noindent\textbf{On ``guaranteed''.} No document can guarantee a result. What this one does is make every recommendation traceable to the professor's own material, and give you a method (Chapter~\ref{ch:method}) that forces you to justify with constraints, numbers and trade-offs, which is what a design question rewards.
\end{abstract}

\tableofcontents

\setcounter{chapter}{-1}
\chapter{How to Attack a Design Question}\label{ch:method}

\section{The six-step answer template}
Every ``design a processor for this scenario'' question can be answered with the same skeleton. Examiners award marks for \emph{linking a constraint to a feature with a reason}, so do this explicitly and in this order.

\begin{enumerate}[leftmargin=*]
\item \textbf{Classify the machine.} Which class from the course: personal mobile device (PMD), desktop, server, cluster/warehouse-scale, supercomputer, embedded? The class fixes what to optimise: PMD = energy and real-time; desktop = price-performance; server = availability, scalability, throughput; embedded = price.
\item \textbf{Classify the workload's parallelism.} Lots of independent requests/threads (TLP/RLP)? Regular loops over arrays (DLP, loop-level ILP)? A single branchy thread (ILP-limited)? Memory-bound with poor locality? This decides core style (wide OoO vs.\ many small cores vs.\ multithreaded).
\item \textbf{Find the bottleneck with a number.} Use the processor-performance equation and AMAT: is $\CPI$ dominated by data hazards, branch penalty, or memory stalls? Quote $f_{branch}\times penalty$, $\AMAT$, miss rate, etc.
\item \textbf{Choose one feature per bottleneck} (from the decision maps in Chapters 2--5), and state \emph{why it matches this workload}.
\item \textbf{State the cost and risk} of each choice (area, power, clock-rate impact, complexity, compatibility). A design with no stated cost looks naive.
\item \textbf{Name what you rejected and why} (``VLIW rejected because binary compatibility is required''). This is where a good answer separates from a good-enough answer.
\end{enumerate}

\begin{key}[One-sentence justification pattern]
``Because the workload is \underline{X} (evidence: \underline{number/constraint}), the bottleneck is \underline{Y}; therefore I choose \underline{feature}, which reduces \underline{CPI term / AMAT term / power} by \underline{mechanism}, at the cost of \underline{cost}; I rejected \underline{alternative} because \underline{constraint}.''
\end{key}

\section{Master lookup: constraint keyword $\rightarrow$ what it implies}
Scan the scenario for these keywords first. The right-hand column is the default direction; later chapters refine it.

\small
\begin{longtable}{L{3.6cm}L{6.3cm}L{6.4cm}}
\toprule
\tblhead{\textbf{If the scenario says\ldots} & \textbf{Lean towards\ldots} & \textbf{Lean away from\ldots}}\\
\midrule
\endfirsthead
\toprule
\tblhead{\textbf{If the scenario says\ldots} & \textbf{Lean towards\ldots} & \textbf{Lean away from\ldots}}\\
\midrule
\endhead
Battery / energy / thermal limit / PMD &
Simple or modestly-wide core, small caches, few cache levels, way-predicted caches, DVFS, power gating, turning off idle cores, multicore of small cores instead of one huge core &
Very wide issue, deep speculation, trace cache, large windows, aggressive prefetch that burns bandwidth (ILP techniques are \emph{energy-inefficient}: slide ``Limits to ILP'') \\
\addlinespace
Low cost / embedded / price-sensitive &
Simple pipeline, static branch prediction / delay slot, direct-mapped small cache, VLIW if code is fixed and compiler is good (simple hardware) &
Tomasulo+ROB, SMT, large L2/L3, large BTB \\
\addlinespace
Server, throughput, many independent requests, database, web &
Multicore + multithreading (fine-grained or SMT), large shared L2/L3, interleaved/banked memory, non-blocking caches, dependability (availability) &
Spending area on single-thread ILP (diminishing returns) \\
\addlinespace
Single-thread latency / desktop / gaming / one big program &
Wide dynamically-scheduled speculative superscalar (Tomasulo+ROB), good branch predictor (two-level / gshare + BTB + return-address stack), small fast L1 + big L2, non-blocking caches, prefetch &
Coarse-grained multithreading (does nothing for one thread), VLIW (if code is irregular) \\
\addlinespace
Regular FP loops, arrays, scientific, matrix &
Loop unrolling + software pipelining + compiler scheduling, pipelined FP units, blocking/interchange/fusion, hardware or compiler prefetch, interleaved memory, VLIW is viable; (vector/SIMD/GPU\bs{}) &
Large dependence on branch prediction (loops are predictable anyway), big victim caches (conflicts are fixable in software) \\
\addlinespace
Binary compatibility across generations / legacy software &
Dynamically scheduled superscalar (hardware finds parallelism; same code runs on next generation) &
VLIW (breaks when number of units or latencies change) \\
\addlinespace
Unpredictable latencies (cache misses) / data-dependent latencies &
Dynamic scheduling (scoreboard/Tomasulo), non-blocking cache, multithreading to hide stalls &
Static scheduling relying on fixed latencies; VLIW lock-step (one stall stalls everything) \\
\addlinespace
Deep pipeline / very high clock &
Strong branch prediction (misprediction cost grows with depth), BTB, small simple L1 so hit time fits cycle, pipelined cache access &
Delayed branch (slots become hard to fill), large associative L1 (hit time $\uparrow$) \\
\addlinespace
High branch frequency (compilers, interpreters, control code) &
Dynamic predictor: 2-bit counters $\rightarrow$ correlating / two-level / gshare; BTB; return-address predictor &
Predict-not-taken alone, 1-bit predictor \\
\addlinespace
Many outstanding misses possible; out-of-order core present &
Non-blocking cache (hit-under-miss / miss-under-miss), multiple memory banks &
Blocking cache (wastes the OoO core) \\
\addlinespace
Large working set, poor locality, pointer-chasing &
Multithreading to hide latency (fine-grained/SMT), large last-level cache, victim cache, deeper hierarchy &
Blocking/tiling and stream prefetch (need regular access patterns) \\
\addlinespace
Memory bandwidth limited, sequential accesses &
Interleaved or wide main memory, larger blocks (up to the AMAT optimum), prefetch stream buffers &
Very large blocks in a small cache \\
\addlinespace
Limited die area &
Multicore of smaller cores (each core simpler, easier to design/manufacture), small-area predictors, SMT as cheap TLP (Power5: SMT cost $\approx$24\% core area) &
Wide issue (logic and ports grow faster than issue width) \\
\addlinespace
Real-time / predictability &
Static techniques (fixed delay slots, static prediction, static schedule); small caches with known behaviour\bs{} &
Heavy speculation and large caches with variable timing\bs{} \\
\addlinespace
High availability (24/7) &
Redundancy, ECC memory, fast repair; compute availability = MTTF/(MTTF+MTTR) &
--- \\
\bottomrule
\end{longtable}
\normalsize

\begin{trap}[Common failure mode]
Choosing ``the most advanced thing'' for every box. The exam rewards \emph{fit}. A PMD answer containing a 12-wide speculative core and a trace cache loses marks; so does a server answer built around a single fast core with no multithreading.
\end{trap}

\section{Roadmap of the chapters}
\begin{center}
\small
\begin{tabular}{L{4.2cm}L{3.2cm}L{8cm}}
\toprule
\tblhead{\textbf{Part of the architecture} & \textbf{Chapter} & \textbf{What you decide there}}\\
\midrule
Fundamentals, power, cost & Ch.~1 & Machine class, power/energy strategy, cost and dependability, Amdahl \\
Pipeline \& branch handling & Ch.~2 & Pipeline depth, FP unit latencies, static vs.\ dynamic prediction, BHT/BTB/gshare \\
Instruction-level parallelism & Ch.~3 & Compiler scheduling, scoreboard vs.\ Tomasulo vs.\ ROB, superscalar vs.\ VLIW, issue width \\
Memory hierarchy & Ch.~4 & Cache organisation, miss/penalty/hit-time techniques, main memory, virtual memory \\
Thread-level parallelism & Ch.~5 & Fine/coarse multithreading, SMT, multicore, shared vs.\ private caches \\
Integration & Ch.~6 & Worked end-to-end scenarios, traps, final checklist \\
\bottomrule
\end{tabular}
\end{center}

\chapter{Fundamentals, Constraints and Trade-offs}\label{ch:fund}
This chapter holds the \emph{constraints} an exam scenario will give you (power, cost, availability, class of machine) and the quantitative tools to reason about them. Everything later is a way of improving one term of the equations here.

\section{Classes of computers: what each one optimises}
\small
\begin{tabularx}{\textwidth}{L{2.8cm}L{4.2cm}X}
\toprule
\tblhead{\textbf{Class} & \textbf{Emphasis (slides)} & \textbf{Design consequence}}\\
\midrule
Personal mobile device (phones, tablets) & Energy efficiency, real-time & Low-power cores, small/few cache levels, DVFS and power gating; throughput per joule beats peak speed. \\
Desktop & Price-performance & Single-thread performance still matters; balanced OoO core with moderate multicore. \\
Server & Availability, scalability, throughput & Multicore/multithreading, large caches, high memory bandwidth, ECC/redundancy. \\
Cluster / warehouse-scale (SaaS) & Availability, price-performance & Request-level parallelism; cost and energy per request dominate. \\
Supercomputer (sub-class) & FP performance, fast internal network & Wide FP units, high memory bandwidth, data-parallel structure. \\
Embedded & Price & Cheapest design that just meets the requirement; simple pipeline, small caches. \\
\bottomrule
\end{tabularx}
\normalsize
\src{ca01 ``Classes of Computers''}

\section{Kinds of parallelism and which hardware exploits them}
\begin{itemize}
\item \textbf{Application-level:} data-level parallelism (DLP: same operation on lots of data) and task-level parallelism (TLP: independent tasks).
\item \textbf{Architecture-level:} instruction-level parallelism (ILP: pipelining, superscalar, VLIW), vector/GPU (DLP), thread-level parallelism (multithreading, multicore), request-level parallelism (clusters/WSC).
\item \textbf{Flynn's taxonomy:} SISD; SIMD (vector architectures, multimedia extensions, GPUs); MISD (no commercial implementation); MIMD (tightly-coupled = multicore/shared memory; loosely-coupled = clusters).
\end{itemize}

\begin{dec}[Match workload to parallelism]
\emph{Same operation on huge regular arrays} $\rightarrow$ SIMD/vector/GPU (the slides only mention these as ``data-level parallelism''; modern graphics cards are SIMD). \emph{Independent threads/requests} $\rightarrow$ MIMD: multicore + multithreading. \emph{One sequential thread} $\rightarrow$ ILP hardware (pipeline, superscalar, speculation), but remember ILP saturates at roughly 3--6 issue.
\end{dec}
\src{ca01 ``Parallelism'', ``Flynn's Taxonomy''; ca08 ``General context: Multiprocessors''}

\begin{trap}[Scope warning]
The course outline lists vector/SIMD/GPU, directory-based coherence, memory consistency and warehouse-scale computers, but \emph{none of the provided slides cover them beyond a definition}. If the exam asks about them, you must supply knowledge beyond this guide.
\end{trap}

\section{The quantitative toolbox}
\subsection{Performance}
\[
\text{CPU time}=\text{IC}\times\CPI\times T_{clk},\qquad
\CPI=\frac{\sum_i \text{IC}_i\cdot\CPI_i}{\text{IC}}
\]
\begin{itemize}
\item \textbf{IC} (instruction count): ISA design and the compiler. \textbf{CPI}: microarchitecture (hazards, branches, memory stalls, issue width). \textbf{$T_{clk}$}: technology and pipeline depth.
\item Every technique in this manual attacks one of these factors, and usually \emph{trades} one against another (deeper pipeline: $T_{clk}\downarrow$ but $\CPI\uparrow$).
\item Speedup of X relative to Y $=$ execution time$_Y$ / execution time$_X$. Benchmarks: kernels, toy programs, synthetic, suites (SPEC, TPC-C).
\end{itemize}
\textbf{Amdahl's law:}\bs\ if a fraction $f$ of execution time is sped up by $s$, overall speedup $=1/\big((1-f)+f/s\big)$. Slide principle: \emph{focus on the common case}.
\src{ca01 ``Principles of Computer Design'', ``Measuring Performance''}

\subsection{Power and energy}
\begin{center}
\begin{tabular}{L{6.2cm}L{9.5cm}}
\toprule
\tblhead{\textbf{Quantity} & \textbf{Formula / statement (slides)}}\\
\midrule
Dynamic energy per switch & $\tfrac12\times C_{load}\times V^2$ \\
Dynamic power & $\tfrac12\times C_{load}\times V^2\times f_{switched}$ \\
Static power & $I_{static}\times V$; scales with number of transistors; reduce by \emph{power gating} \\
TDP (thermal design power) & Sustained power target for supply and cooling; lower than peak, higher than average \\
Better metric than power & Energy per task \\
\bottomrule
\end{tabular}
\end{center}
\begin{itemize}
\item \textbf{Reducing the clock rate reduces power but not energy per task} (slide). Lowering \emph{voltage} (DVFS) reduces energy because of the $V^2$ term.
\item Techniques the slides list: dynamic voltage-frequency scaling, low-power states for DRAM/disks, turning off cores, overclocking within thermal limits, power gating (static power).
\item Scale of the problem: Intel 80386 $\approx$2\,W; 3.3\,GHz Core i7 $\approx$130\,W from a 1.5\,cm$\times$1.5\,cm die, ``the limit of what can be cooled by air''.
\item If voltage scales with frequency, dynamic power falls roughly with $f^3$ (0.8$\times$ clock $\Rightarrow$ $\approx$0.51$\times$ power)\bs; this is why \emph{two slower cores can beat one fast core} on throughput per watt.
\end{itemize}

\begin{pick}[Energy-first design]
Choose DVFS/power gating, small caches, narrow issue and multicore of small cores whenever the scenario says battery, thermal envelope, or energy cost per request. Justify with $P\propto CV^2f$ and ``energy per task is the right metric''.
\end{pick}
\src{ca01 ``Power and Energy'', ``Dynamic Energy and Power'', ``Reducing Power'', ``Static Power''; ca06 ``Limits to ILP''}

\subsection{Cost}
\begin{itemize}
\item Cost falls along a \emph{learning curve} (yield improves). DRAM price tracks cost; microprocessor price depends on volume (10\% less per doubling of volume).
\item Die yield follows the Bose--Einstein form $\text{yield}=\text{wafer yield}\times(1+\text{defects/area}\times\text{die area})^{-N}$ with $N$ the process-complexity factor (slide: 11.5--15.5; defects 0.016--0.057/cm$^2$)\bs\ (formula is an image in the slide). \textbf{Consequence: bigger dies are disproportionately more expensive}, favouring several small cores over one huge core.
\end{itemize}
\src{ca01 ``Trends in Cost'', ``Integrated Circuit Cost''}

\subsection{Dependability}
\[
\text{MTBF}=\text{MTTF}+\text{MTTR},\qquad \text{Availability}=\frac{\text{MTTF}}{\text{MTBF}}
\]
For servers (availability is their emphasis) quote this and add redundancy and ECC (the slides mention error-correcting codes when discussing wide memory). \src{ca01 ``Dependability''; ca05 ``Wider Main Memory''}

\section{Technology trends that justify design choices}
\small
\begin{tabularx}{\textwidth}{L{4.2cm}X}
\toprule
\tblhead{\textbf{Trend (slide numbers)} & \textbf{What it implies for design}}\\
\midrule
Bandwidth improved 10{,}000--25{,}000$\times$ for processors but only 300--1200$\times$ for memory/disks; latency improved 30--80$\times$ vs.\ 6--8$\times$ &
Memory is the bottleneck (``memory wall''). Use caches, interleaving/wide memory, prefetching, and latency-hiding (OoO, non-blocking caches, multithreading). Bandwidth is easy to buy; latency is not. \\
Transistor density $+35\%$/yr; transistor speed scales, \textbf{wire delay does not} &
Huge centralised structures (many register ports, big CDB, wide issue logic) hurt clock rate. Favours smaller replicated cores and simpler structures. \\
Single-processor performance improvement ended in 2003; ILP cannot be leveraged further &
Move to explicit parallelism: DLP, TLP, RLP (``require explicit restructuring of the application''). Justifies multicore. \\
DRAM capacity growth slowing (25--40\%/yr); flash 15--20$\times$ cheaper per bit than DRAM; disk 15--25$\times$ cheaper than flash &
Deeper hierarchy: SRAM caches $\rightarrow$ DRAM $\rightarrow$ flash $\rightarrow$ disk, each level sized by cost per bit. \\
\bottomrule
\end{tabularx}
\normalsize
\src{ca01 ``Trends in Technology'', ``Bandwidth and Latency'', ``Transistors and Wires'', ``Current Trends in Architecture''}

\section{Chapter 1 cheat-sheet}
\begin{key}
\begin{itemize}
\item Optimise \emph{energy per task}, not clock rate, for mobile/server cost; lower voltage is the real lever.
\item Bigger die $\Rightarrow$ worse yield $\Rightarrow$ cost grows faster than area.
\item Bandwidth grows much faster than latency: spend hardware on hiding latency.
\item ILP is capped (3--6 issue); explicit TLP/DLP is the next step.
\item Availability $=$ MTTF$/($MTTF$+$MTTR$)$.
\end{itemize}
\end{key}

\chapter{Pipeline, Long-Latency Units and Branch Handling}\label{ch:pipe}

\section{Decision map}
\small
\begin{tabularx}{\textwidth}{L{4.0cm}L{4.4cm}X}
\toprule
\tblhead{\textbf{Problem observed} & \textbf{Feature to apply} & \textbf{Why it fits}}\\
\midrule
Stalls from RAW hazards between dependent ALU ops & Forwarding/bypassing & Result is available at end of EX; forward instead of waiting for write-back. \\
Load followed by use & Compiler load-delay scheduling; forwarding from MEM & Load data is late; fill the slot(s) with independent instructions. \\
Control hazard, shallow pipe, predictable branches & Early PC evaluation in ID + static prediction / delay slot & Cheapest (no predictor memory), penalty already only 1 cycle. \\
Control hazard, deep pipe or high branch frequency & Dynamic prediction (BHT $\rightarrow$ correlating/two-level $\rightarrow$ gshare) + BTB & Penalty $=$ several cycles; accuracy matters more than anything. \\
Loop-heavy code & BTFNT or 2-bit counter & Back-edge is taken almost always; 2-bit avoids the double misprediction. \\
Target address needed early & BTB (and target buffer / branch folding) & Fetch target in IF with 0--1 bubble. \\
Returns / indirect jumps & Return-address stack & Targets vary at run time; stack predicts them. \\
Divide-heavy or multiply-heavy FP code & Pipelined (or replicated) FP units & Initiation interval 25 for an unpipelined divider blocks every other divide. \\
Want higher clock & Deeper pipeline (R4000 8-stage) & Shorter stages; but raises load/branch delay. \\
\bottomrule
\end{tabularx}
\normalsize

\section{Pipeline hazards: the three kinds and their standard cures}
\begin{center}
\small
\begin{tabular}{L{2.6cm}L{4.8cm}L{8cm}}
\toprule
\tblhead{\textbf{Hazard} & \textbf{Cause} & \textbf{Cure (slides)}}\\
\midrule
Structural & Two instructions need the same hardware in the same cycle & More hardware resources (replicate/pipeline units, more ports). \\
Data (RAW/WAR/WAW) & Dependence between instructions & Forwarding; compiler scheduling; (hardware) register renaming for WAR/WAW. \\
Control & Next PC unknown until branch resolves & Early evaluation of PC, delayed branch, prediction. \\
\bottomrule
\end{tabular}
\end{center}
\begin{itemize}
\item Pipelining improves instruction \emph{bandwidth} (throughput), not the latency of one instruction.
\item \textbf{Longer pipelines increase the impact of hazards}; interrupts, complex ISAs and FP all make pipelining harder.
\item Compilers reduce the cost of data and control hazards: load delay slots, branch delay slots, branch prediction hints.
\end{itemize}
\src{ca03\_noor ``Review: Summary of Pipelining Basics''}

\section{Long-latency functional units (FP pipeline)}
The slides' running example is a RISC-V/MIPS-style pipeline in which EX may repeat, and there are several functional units.

\begin{center}
\begin{tabular}{L{6cm}C{3cm}C{3.6cm}}
\toprule
\tblhead{\textbf{Functional unit} & \textbf{Latency} & \textbf{Initiation interval}}\\
\midrule
Integer ALU & 0 & 1 \\
Data memory (integer and FP loads) & 1 & 1 \\
FP add (4-stage pipelined) & 3 & 1 \\
FP/integer multiply (7-stage pipelined) & 6 & 1 \\
FP/integer divide (not pipelined) & 24 & 25 \\
\bottomrule
\end{tabular}
\end{center}
\textbf{Latency} $=$ intervening cycles between the producer and a consumer that avoids a RAW stall. \textbf{Initiation interval} $=$ cycles that must separate two operations of the \emph{same} type issued to the unit.

Latency table used in the compiler-scheduling examples (cycles between producer and consumer): FP ALU$\rightarrow$FP ALU 3; FP ALU$\rightarrow$store 2; load$\rightarrow$FP ALU 1; load$\rightarrow$store 0; integer$\rightarrow$integer 0.

\begin{itemize}
\item Longer pipeline $\Rightarrow$ more frequent RAW stalls (e.g.\ the 4th instruction after an FP add is the first that can use its result without a stall).
\item New hazards appear: \textbf{structural hazard at the register-file write port} (three instructions wanting write-back in the same cycle), \textbf{structural hazard on the unpipelined divider}, and \textbf{WAW hazards / out-of-order completion} because a short operation can finish before an earlier long one.
\end{itemize}

\begin{pick}[Pipelined vs.\ unpipelined divider]
Pipelined divider (or several dividers) when the code has frequent divides close together: with initiation interval 25, a second divide stalls for up to 24 cycles. An unpipelined divider is a cheap, acceptable choice when divides are rare, since it needs far less area and power. Adder and multiplier are worth pipelining almost always (initiation interval 1).
\end{pick}
\src{ca03\_noorNew ``RISC V FP Operation'', ``Latencies and initiation intervals'', ``A pipeline that supports multiple outstanding FP operations''; ca03\_noor ``FP Loop: Where are the Hazards?''}

\section{Deeper pipelines: the MIPS R4000 case study}
The R4000 splits the 5-stage integer pipeline into \textbf{8 stages}: IF, IS, RF, EX, DF, DS, TC, WB (cache access takes two stages each for instruction and data; TC is the tag check). The aim is a \textbf{higher clock rate}.

\begin{center}
\small
\begin{tabular}{L{5cm}L{10cm}}
\toprule
\tblhead{\textbf{Effect} & \textbf{Value / reason}}\\
\midrule
Load delay & \textbf{2 cycles} (the older deck states ``TWO cycle load latency''; the newer deck's slide title reads ``x1'' but its pipeline table shows two Stall cycles --- use 2). Data is available at end of DS and is bypassed; if TC reports a miss the pipeline backs up. \\
Branch delay & \textbf{3 cycles}: condition evaluated in EX; one delay slot plus two stalls; ``branch likely'' cancels the delay slot when not taken. \\
FP unit & Three units: divider, multiplier, adder; adder also used in the last step of multiply and divide. Stage kinds: U (unpack), S (operand shift), A (mantissa add), R (round), E (exception test), M, N (multiplier stages), D (divide stage, used 28 times in a row for divide). \\
\bottomrule
\end{tabular}
\end{center}

\begin{dec}[Deep vs.\ shallow pipeline]
\emph{Go deeper} when the target is a high clock rate and the code has good branch predictability and few back-to-back load-use pairs (scientific/loop code), \emph{and} you can afford a strong branch predictor. \emph{Stay shallow} for branchy, low-ILP, low-power or area-limited designs: every extra stage adds to the misprediction and load penalty (and, in general, to power\bs) (deeper pipelines also magnify branch cost: slide ``Branch Prediction Techniques'').
\end{dec}
\src{ca03\_noor ``Case Study: MIPS R4000''; ca03\_noorNew ``The MIPS R4000 Pipeline'' through ``The Floating-Point Pipeline''}

\section{The control hazard problem in numbers}
\subsection{Where the branch resolves determines the penalty (5-stage MIPS)}
\begin{center}
\begin{tabular}{L{8.6cm}C{3.2cm}}
\toprule
\tblhead{\textbf{Scheme} & \textbf{Stall cycles per branch}}\\
\midrule
Resolve at MEM, no forwarding (conservative stall) & 3 \\
Resolve at EX (with forwarding) & 2 \\
\textbf{Early evaluation: compare registers, compute target and update PC in ID} & \textbf{1} \\
\bottomrule
\end{tabular}
\end{center}
Early evaluation has a side-effect: an ALU instruction that produces the branch operand needs \textbf{1 extra stall}; a load that produces it needs \textbf{2 extra stalls}.

\[
\text{stall cycles per instruction (branches)} = f_{branch}\times \text{penalty};\qquad
\CPI = \CPI_{base}+f_{branch}\times(\text{mispredict rate})\times\text{penalty}\bs
\]
Even a 1-cycle penalty costs 10--30\% performance depending on branch frequency (slide). \textbf{Importance of prediction grows with (a) pipeline depth (penalty) and (b) branch frequency.}

\begin{ex}[Why prediction matters more in deep pipes]\bs
$f_{branch}=20\%$. Shallow pipe with early evaluation: always stall 1 cycle $\Rightarrow$ $+0.20$ CPI. Deep pipe, penalty 3: always stall $\Rightarrow +0.60$ CPI; with a 90\%-accurate predictor $\Rightarrow 0.20\times0.10\times3=+0.06$ CPI. The predictor recovers 90\% of the loss.
\end{ex}
\src{ca04 ``The Problem of Control Hazards'' through ``Branch Prediction Techniques''}

\section{Static branch prediction (decided at compile time)}
Static schemes need no run-time hardware, so they are used where branch behaviour is highly predictable at compile time, and also to assist dynamic predictors.

\small
\begin{longtable}{L{2.9cm}L{4.4cm}L{4.5cm}L{3.9cm}}
\toprule
\tblhead{\textbf{Scheme} & \textbf{How it works} & \textbf{Best scenario} & \textbf{Weakness}}\\
\midrule
\endhead
Predict not taken & Keep fetching fall-through; flush 1 instruction if taken (1-cycle penalty) & Branches mostly not taken (error checks, forward branches); simplest hardware. If 50\% not taken, halves control-hazard cost & Loops (back-edges are taken) \\
Predict taken & Fetch target once decoded & Only when the \emph{target is known before the outcome} (needs BTA in IF or a BTB). No advantage in the base MIPS pipeline & Needs extra target hardware \\
BTFNT (backward taken, forward not taken) & Direction of branch decides & Loop-dominated code: back-edge branches almost always taken & Irregular forward branches \\
Profile-driven & Compiler uses profile from earlier runs/hints & Stable behaviour across inputs (fixed embedded workloads) & Behaviour that changes with input \\
Delayed branch & Compiler fills slot(s) after branch with a useful instruction; executed whether or not taken & Shallow pipe with a single slot; compiler fills $\approx$50\% usefully & Deep pipelines (several slots, hard to fill); remaining slots become NOPs \\
\bottomrule
\end{longtable}
\normalsize

\subsection{Choosing where the delay-slot instruction comes from}
\begin{center}
\small
\begin{tabular}{L{2.6cm}L{6cm}L{6.2cm}}
\toprule
\tblhead{\textbf{Source} & \textbf{Use when} & \textbf{Legality}}\\
\midrule
From before & Always valid if an independent instruction exists before the branch & Safe: executes the same either way \\
From target & \textbf{Branch usually taken} (loop back-edges); target instruction must be copied & Must be harmless (e.g.\ writes an unused temporary) if branch falls through \\
From fall-through & \textbf{Branch usually not taken} (forward branches) & Must be harmless if branch is taken \\
\bottomrule
\end{tabular}
\end{center}
\textbf{Canceling (nullifying) branch:} the instruction encodes the predicted direction; if the prediction is wrong the delay-slot instruction becomes a NOP, so the compiler need not be conservative. MIPS \emph{branch-likely} is cancel-if-not-taken (slot filled from target) --- ideal for loop back-edges.
\src{ca04 ``Static Branch Prediction Techniques'' through ``Delayed Branch Technique'' (all pages)}

\section{Dynamic branch prediction (decided by hardware at run time)}
Two cooperating pieces: a \textbf{branch-outcome predictor} (taken?) and a \textbf{branch-target predictor/BTB} (where to?). The fetch unit uses them to pick the next PC in IF.

\subsection{Branch history table (BHT) and counters}
\begin{itemize}
\item \textbf{1-bit BHT:} table indexed by low-order bits of the branch PC; no tags (every access ``hits''; aliasing is tolerated because the entry is only a hint). Weakness: a loop branch (taken 9, not-taken 1) mispredicts \emph{twice} per loop execution (at exit, and on re-entry) $\Rightarrow$ only \textbf{80\% accuracy}.
\item \textbf{2-bit saturating counter:} must be wrong twice before the prediction flips; loop exit costs only one misprediction. Predict taken if counter $\ge$ half of max.
\item \textbf{n-bit counter:} generalisation; studies show 2-bit performs almost as well, so \emph{use 2-bit}.
\item A 4K-entry 2-bit BHT achieves 82--99\% accuracy (1--18\% misprediction) on SPEC89 and is almost as good as an infinite table; so beyond $\sim$4K entries, make the \emph{predictor smarter}, not bigger. Aliasing is cured by more rows or a hash.
\end{itemize}

\subsection{Correlating (two-level) predictors}
Idea: the outcome of recent \emph{other} branches helps predict this one (e.g.\ if the two preceding branches were both not taken, branch bb3 will be taken).
\begin{itemize}
\item \textbf{$(m,n)$ predictor:} the last $m$ branch outcomes (global history) select one of $2^m$ BHTs, each with $n$-bit counters. $(0,2)$ is the ordinary 2-bit predictor. $(2,2)$: four 2-bit BHTs selected by 2-bit history.
\item Index $=$ low-order PC bits concatenated with the $m$-bit history. Example: a $(2,2)$ predictor with 64 total entries uses a 6-bit index $=$ 2 history bits + 4 PC bits; each BHT has 16 entries.
\item A small $(2,2)$ predictor beats a plain 2-bit predictor of the same storage and often beats an unlimited-size 2-bit predictor.
\item \textbf{Two-level adaptive:} first level $=$ branch history register (BHR, a $k$-bit shift register of recent outcomes); second level $=$ pattern history table (PHT) of 2-bit counters indexed by the BHR.
\item \textbf{GAs:} PHT indexed by PC bits \emph{and} global history (local + global). \textbf{gshare:} PHT indexed by \textbf{XOR} of low PC bits with the BHR (local XOR global), which uses the table better than concatenation\bs.
\end{itemize}

\subsection{Branch target buffer (BTB) and fetch-unit support}
\begin{itemize}
\item BTB is a \emph{cache} accessed in IF with the PC of the instruction being fetched. Entry: \textbf{tag + predicted target (PC-relative) + optional prediction bits}. Only \textbf{taken} branches need an entry. The tag \emph{must} match (a wrong target cannot be used, unlike BHT aliasing).
\item If BTB hit and predicted taken: fetch from the predicted PC immediately. On misprediction: kill fetched instruction, restart at the other target, update/delete the entry. Correct prediction: no stalls.
\item \textbf{Target instruction buffers} store the target \emph{instruction}(s): faster, allow a larger BTB, and permit \textbf{branch folding} (0-cycle unconditional branches).
\item \textbf{Integrated instruction-fetch unit:} autonomous unit doing branch prediction, prefetch and instruction buffering to feed a wide back end.
\item \textbf{Return-address predictor:} small stack of recent return addresses, for indirect jumps/returns whose destination varies at run time.
\end{itemize}

\subsection{Which dynamic predictor for which scenario}
\small
\begin{longtable}{L{3.7cm}L{4.4cm}L{7.3cm}}
\toprule
\tblhead{\textbf{Scenario} & \textbf{Choose} & \textbf{Reason}}\\
\midrule
\endhead
Small/cheap core, loop code & 2-bit BHT & Fixes the loop double-miss at tiny cost (2 bits/entry, no tags). \\
Control-heavy integer code (compilers, interpreters), correlated if-chains & Correlating $(2,2)$ or two-level (GAs/gshare) & Uses other branches' outcomes; small table beats a much larger 2-bit table. \\
Same storage budget, want best accuracy & gshare & XOR indexing spreads history and PC across the table. \\
Wide-issue or deep pipeline & Two-level predictor + BTB + return-address stack, feeding \emph{speculation} & Cost of a miss is large; wide issue may need a branch every cycle. \\
Many function calls/returns & Return-address stack & Returns are indirect with run-time targets. \\
Fixed embedded workload, no predictor area & Static: profile-driven or BTFNT (+delay slot) & No hardware; behaviour known offline. \\
Branch target needed in IF for 0--1 bubble taken branches & BTB (+ target buffer for folding) & Predicts the \emph{address}, not just direction. \\
\bottomrule
\end{longtable}
\normalsize

\begin{trap}[Branch-prediction traps]
\begin{itemize}
\item 1-bit loop example: accuracy is 80\% for a 10-iteration loop (8 correct, 2 wrong), not 90\%.
\item ``Predict taken'' is useless unless the target address is available before the outcome (needs a BTB or early BTA).
\item A BHT has no tags (hint only); a BTB \emph{does} (an incorrect address cannot be used).
\item Prediction alone is not enough for wide-issue machines: it must be combined with \textbf{speculation} (Chapter~\ref{ch:ilp}).
\end{itemize}
\end{trap}
\src{ca04 ``Dynamic Branch Prediction Techniques'' through ``Branch Target Buffer'', ``High Performance Instruction Delivery''}

\section{Chapter 2 cheat-sheet}
\begin{key}
\begin{itemize}
\item Penalty: MEM 3 $\rightarrow$ EX 2 $\rightarrow$ ID (early eval.) 1. Stalls/instr $=f_{branch}\times$ penalty.
\item Deeper pipe $\Rightarrow$ bigger penalty $\Rightarrow$ dynamic prediction (two-level + BTB + RAS).
\item 1-bit: 2 misses per loop; 2-bit: 1 miss. $(m,n)$: $2^m$ tables of $n$-bit counters. gshare: PC XOR history.
\item Delay slot from target for loops, from fall-through for forward branches; $\approx$50\% usefully filled.
\item FP latencies 3 / 6 / 24 (initiation 1 / 1 / 25). R4000: load delay 2, branch delay 3.
\end{itemize}
\end{key}

\chapter{Instruction-Level Parallelism: Scheduling, Speculation, Multiple Issue}\label{ch:ilp}

\section{Decision map}
\small
\begin{tabularx}{\textwidth}{L{4.2cm}L{4.3cm}X}
\toprule
\tblhead{\textbf{Scenario} & \textbf{Feature} & \textbf{Why}}\\
\midrule
Regular loops, independent iterations, fixed known latencies & Loop unrolling + static scheduling (+ software pipelining) & Compiler creates independent instructions to fill latency slots and removes loop overhead. \\
Unknown dependences at compile time (memory aliasing, cache-miss latencies) & Dynamic scheduling (Tomasulo / scoreboard) & Hardware reorders at run time; compiler need not know the microarchitecture. \\
Few architectural registers, lots of WAR/WAW & Tomasulo (register renaming in reservation stations) & Renames in hardware to a larger virtual register set. \\
Simplest possible out-of-order, few non-pipelined FUs & Scoreboard (CDC 6600 style) & Cheap centralised control; accepts stalls on WAR/WAW and no forwarding. \\
Need precise exceptions + branch speculation & Tomasulo $+$ reorder buffer (ROB) & Out-of-order execute, in-order commit; easy undo. \\
Need CPI $<1$ & Multiple issue: superscalar or VLIW & More than one instruction per cycle. \\
Binary compatibility across generations & Dynamically scheduled superscalar & Same binary, hardware finds ILP. \\
Fixed-product, compiler-controlled, power/area simplicity & VLIW (static scheduling) & No issue/hazard-detection hardware. \\
Wider than $\sim$4--6 issue is requested & Don't: go multicore/TLP instead & ILP saturates; energy inefficient; clock rate suffers. \\
\bottomrule
\end{tabularx}
\normalsize

\section{Dependences: what limits ILP}
\begin{center}
\small
\begin{tabular}{L{3.2cm}L{5.2cm}L{6.8cm}}
\toprule
\tblhead{\textbf{Type} & \textbf{Meaning} & \textbf{Can it be removed?}}\\
\midrule
True (data) dependence & Instruction $j$ uses a result produced by $i$ (RAW if a hazard) & \textbf{No.} It is the real data flow; instructions cannot execute in parallel. \\
Anti-dependence (name) & $j$ writes a register/memory location that earlier $i$ reads (WAR) & Yes: \emph{register renaming} (just a naming conflict). \\
Output dependence (name) & $i$ and $j$ write the same location (WAW); order must be kept & Yes: register renaming. \\
Control dependence & Instruction executes only if a branch goes a certain way & Relaxed by speculation/prediction; must preserve exception behaviour and data flow. \\
\bottomrule
\end{tabular}
\end{center}
\begin{itemize}
\item Register dependences are easy to detect (fixed names). \textbf{Memory dependences are hard}: does \texttt{100(R4)} equal \texttt{20(R6)}? Do different iterations touch the same address? (alias analysis).
\item \textbf{Control rules for code motion:} an instruction control-dependent on a branch cannot be moved before it; an instruction not control-dependent on it cannot be moved after it.
\item \textbf{Loop-carried dependence:} e.g.\ \texttt{A[i+1]=A[i]+C[i]} depends on the previous iteration, so iterations cannot run in parallel and the loop cannot be freely unrolled/parallelised. Safe to unroll when iterations are independent.
\end{itemize}
\src{ca03\_noor ``Compiler Perspectives on Code Movement'', ``When Safe to Unroll Loop?''; lec06 ``Review: Exceptions and Compiler Scheduling''}

\section{Static (compiler) scheduling}
\subsection{Loop unrolling: the slides' worked example}
Loop: \texttt{LD F0,0(R1); ADDD F4,F0,F2; SD 0(R1),F4; SUBI R1,R1,8; BNEZ R1,Loop} (+ delay slot).

\begin{center}
\small
\begin{tabular}{L{8.5cm}C{2.8cm}C{3.4cm}}
\toprule
\tblhead{\textbf{Version} & \textbf{Cycles} & \textbf{Cycles / iteration}}\\
\midrule
Original (stalls after LD, ADDD, and in branch slot) & 9 & 9 \\
Reordered: SUBI moved up, store address adjusted (\texttt{SD 8(R1),F4}) to fill delay slot & 6 & 6 \\
Unrolled $\times$4, straightforward (still stalls) & $15+4\times(1+2)=27$ & 6.75 \\
Unrolled $\times$4, scheduled with renamed registers (loads together, ADDDs together, stores together) & 14 & \textbf{3.5} \\
\bottomrule
\end{tabular}
\end{center}
\textbf{Why it works:} loop overhead (SUBI, BNEZ) is amortised; independent instructions from different iterations fill the 1-, 2- and 3-cycle latency gaps. \textbf{What it requires:} (a) \emph{renaming} (use F0/F4, F6/F8, \dots\ instead of reusing F0/F4, otherwise name dependences stall you), (b) knowledge that \texttt{0(R1)}, \texttt{-8(R1)}, \dots\ are different addresses so loads can move before stores, (c) no loop-carried dependence.

\subsection{Software pipelining and static renaming}
\begin{itemize}
\item \textbf{Loop unrolling:} multiple iterations in the loop body, all instructions from all iterations; needs compiler register renaming.
\item \textbf{Software pipelining:} one instruction from each of several iterations in the body; turns dependences into anti-dependences; often used for \emph{floating-point} with long latencies.
\item Static renaming limits: dependence is only resolved within a window, and \textbf{increases register pressure} for the compiler.
\end{itemize}

\begin{pick}[Static scheduling / unrolling]
Regular numeric loops with independent iterations and \emph{known fixed latencies}; in-order or VLIW hardware; embedded/DSP-style kernels where simple hardware and predictable timing matter. Best when many registers are available.
\end{pick}
\begin{avoid}
Code with unknown memory aliasing, pointer chasing, data-dependent branches, or variable latency from cache misses; or when the same binary must run efficiently on future hardware with different latencies.
\end{avoid}
\bigskip
\src{ca03\_noor ``FP Loop Hazards'' through ``Unrolled Loop That Minimizes Stalls''; lec06 ``Review''}

\section{Dynamic scheduling: why do it in hardware}
\begin{itemize}
\item Works when the dependence cannot be known at compile time; keeps the \emph{compiler simple}; code tuned for one machine runs well on another.
\item Key idea: \textbf{allow instructions behind a stalled one to proceed} (e.g.\ in \texttt{DIVD F0,F2,F4; ADDD F10,F0,F8; SUBD F12,F8,F14}, SUBD is independent and need not wait for the divide).
\item Consequence: out-of-order execution $\Rightarrow$ out-of-order completion $\Rightarrow$ \textbf{WAR and WAW hazards} become possible, and \textbf{precise exceptions} become hard.
\item Disadvantages: substantial hardware complexity; complicates exceptions.
\end{itemize}

\subsection{Scoreboard (CDC 6600, 1963)}
\textbf{In-order issue, out-of-order execution and commit.} Four stages replace ID/EX/WB:
\begin{enumerate}
\item \textbf{Issue:} functional unit free \emph{and} no active instruction with the same destination (no WAW); otherwise stall, and nothing behind issues.
\item \textbf{Read operands:} wait until both operands are available (RAW resolved here); no forwarding.
\item \textbf{Execute.}
\item \textbf{Write result:} stall if a WAR hazard exists with an earlier instruction that has not yet read its operand.
\end{enumerate}
Three tables: instruction status; functional-unit status (Busy, Op, Fi, Fj, Fk, Qj, Qk, Rj, Rk); register result status.

\textbf{Limitations (slide):} no forwarding hardware; limited to instructions in a basic block (small window, $\le5$ instructions in comparison table); few functional units (structural hazards, especially integer/load-store); no issue on structural hazard; \emph{waits} for WAR; \emph{prevents} WAW.

\subsection{Tomasulo's algorithm (IBM 360/91, 1966)}
Goal: high performance \emph{without special compilers}. The 360/91 had only \textbf{4 FP registers}, which blocked compiler scheduling; Tomasulo therefore \textbf{renamed registers in hardware}.
\begin{itemize}
\item \textbf{Distributed control:} \emph{reservation stations (RS)} sit with each functional unit and hold the op, operand values $V_j,V_k$, and tags $Q_j,Q_k$ (the RS that will produce the operand; $Q=0$ means ready). Register specifiers are replaced by values or RS tags $\Rightarrow$ \textbf{register renaming} that \textbf{eliminates WAR and WAW}. More RS than registers $\Rightarrow$ optimisations compilers cannot do.
\item \textbf{Common data bus (CDB):} results broadcast with their \emph{source} tag (``come-from'' bus: 64 data bits + 4 FU-id bits) to all RS, store buffers and the register file. Only the last writer updates the register.
\item Loads and stores are treated as functional units with buffers (load/store buffers hold address and data); memory order is kept through effective-address computation.
\item Three steps: \textbf{Issue} (if RS free, send operands/tags; else stall), \textbf{Execute} (when operands ready; watch the CDB otherwise), \textbf{Write result} (on CDB). No instruction starts until all earlier branches have resolved (non-speculative version).
\item \textbf{Overlaps loop iterations:} renaming gives different iterations different destinations (``dynamic loop unrolling''); RS let integer instructions run ahead past branches so FP ops beyond a basic block can be queued; also hides cache misses.
\end{itemize}

\begin{center}
\small
\begin{tabular}{L{3cm}L{6.2cm}L{6.2cm}}
\toprule
\tblhead{\textbf{Aspect} & \textbf{Tomasulo (IBM 360/91)} & \textbf{Scoreboard (CDC 6600)}}\\
\midrule
Functional units & Pipelined: 6 load, 3 store, 3 add, 2 mult/div & Multiple, not necessarily pipelined: 1 load/store, 1 add, 2 mult, 1 div \\
Window size & $\le 14$ instructions & $\le 5$ instructions \\
Structural hazard & No issue & Same \\
WAR & Renaming avoids stall & Stall completion \\
WAW & Renaming avoids stall & Stall issue \\
Result delivery & Broadcast from FU on CDB & Write/read registers (no forwarding) \\
Control & Distributed (reservation stations) & Centralised scoreboard \\
Worked example (6-instr.\ sequence) & Last instruction writes in cycle 57 & Same code finishes in cycle 62 \\
\bottomrule
\end{tabular}
\end{center}
Scoreboard is slower because of \textbf{structural hazards and lack of forwarding}.

\textbf{Tomasulo drawbacks (slide):} complexity; many associative stores (CDB) at high speed; \textbf{performance limited by the CDB} --- each CDB goes to many units (high capacitance, wiring density), and only \textbf{one FU can complete per cycle per CDB}; more CDBs need more associative logic; \textbf{imprecise interrupts} (fixed by the ROB).

\subsection{Hardware speculation: Tomasulo $+$ reorder buffer (ROB)}
Combines three ideas: \textbf{dynamic branch prediction} (choose what to run), \textbf{speculation} (execute before control dependences resolve, with ability to undo), \textbf{dynamic scheduling} (schedule across basic blocks). Core principle: \textbf{execute out of order, commit in order}; no irrevocable action (register/memory write, exception) before commit.

\begin{itemize}
\item \textbf{ROB:} circular buffer allocated in program order; entries hold instruction type (branch / store / register op), destination (register number or memory address), value, ready bit. It supplies operands between completion and commit and \emph{replaces the store buffer}.
\item \textbf{Four steps:} Issue (needs a free RS \emph{and} ROB slot) $\rightarrow$ Execute $\rightarrow$ Write result (CDB $\rightarrow$ ROB) $\rightarrow$ \textbf{Commit} (at ROB head: write register/memory; a mispredicted branch at the head \textbf{flushes the ROB} and restarts at the correct successor).
\item \textbf{Benefits:} \emph{precise interrupts} (exceptions recorded in the ROB, recognised only at commit; a speculative instruction's exception is flushed with it); \emph{early recovery} from mispredictions (clear entries after the branch, keep older ones); \emph{memory hazards}: stores update memory only at the head, eliminating WAR/WAW through memory; RAW through memory is kept by not letting a load proceed if an earlier store has a matching address.
\item \textbf{Costs:} large control complexity; associative network to find the latest version of a register (or a future file / register-status buffer); as many ROB ports as register-file ports; ROB sized to window.
\item \textbf{Design questions:} register renaming (extended physical registers) \emph{versus} ROB for holding values; \textbf{how much to speculate} --- speculation is hurt by exceptional events (cache miss, TLB miss) and speculating through multiple branches complicates recovery.
\end{itemize}

\begin{dec}[Scoreboard vs.\ Tomasulo vs.\ Tomasulo + ROB]
\begin{itemize}
\item \textbf{No dynamic scheduling} (in-order + compiler): cheapest/lowest power; embedded, predictable loops.
\item \textbf{Scoreboard:} want some out-of-order tolerance with minimal logic and few, non-pipelined units; accept WAR/WAW stalls.
\item \textbf{Tomasulo:} FP-heavy code with long latencies, \emph{few architectural registers}, or variable latencies (cache misses); you want renaming and loop overlap. Cap: one result per CDB per cycle.
\item \textbf{Tomasulo + ROB (speculative):} high-performance desktop/server cores needing precise exceptions and branch speculation, especially with multiple issue. Most area/power.
\end{itemize}
\end{dec}
\src{ca03\_noor (``HW Schemes'', ``Scoreboard'', ``Dynamic Scheduling - Tomasulo''); lec06 (whole deck: ``Tomasulo vs.\ Scoreboard'', ``Tomasulo Drawbacks'', ``Four Steps of Speculative Tomasulo''); ca04 ``Hardware-Based Speculation''}

\section{Multiple issue: superscalar vs.\ VLIW}
To get CPI $<1$ you must complete multiple instructions per clock. Three approaches: \emph{statically scheduled superscalar}, \emph{VLIW}, \emph{dynamically scheduled superscalar}.

\begin{center}
\small
\begin{tabular}{L{3cm}L{6.1cm}L{6.3cm}}
\toprule
\tblhead{\textbf{} & \textbf{Superscalar (dynamic)} & \textbf{VLIW (static)}}\\
\midrule
Who finds parallelism & Hardware at run time (RAW: forwarding; WAR/WAW: renaming) & Compiler at compile time: bundles of $K$ independent operations \\
Code size & Smaller & Larger: NOPs when bundles cannot be filled; unrolled loops \\
Binary compatibility & \textbf{Yes} across hardware generations & \textbf{Weak}: changing number of units ($K\to K+5$) or latencies breaks the schedule; fix = emulator, which hurts performance \\
Hardware & Complex issue/decode, dependence checking, renaming & \textbf{Simple decode/issue}, no interlock hardware (compiler checks) \\
Registers/ports & --- & More registers; simplified; may need multiple register files; VLIW example needs 7 read/3 write int and 5 read/3 write FP ports \\
Stalls & Independent instructions continue & \textbf{Lock-step}: one hazard stalls the whole bundle \\
Examples & Pentium 4, Athlon 64, Power5 & Itanium 2 (IA-64/EPIC) \\
\bottomrule
\end{tabular}
\end{center}
Example VLIW word: one integer op (or branch), two independent FP ops, two independent memory references. There must be enough parallelism to fill the slots.

\textbf{Dynamic multiple issue with speculation} (modern microarchitecture): limit instruction classes per bundle (e.g.\ one FP, one int, one load, one store); examine dependences among the bundle's instructions and encode them in the reservation stations; need multiple completion/commit (multiple CDBs). \textbf{Issue logic becomes the bottleneck.}

\begin{pick}[Superscalar (dynamic)]
General-purpose products with many software generations, legacy binaries, irregular integer code and unpredictable latencies.
\end{pick}
\begin{pick}[VLIW]
Embedded/DSP or fixed-function products where the compiler is controlled, the code is regular (loops with abundant ILP), the hardware must be simple, and binary compatibility across generations is \emph{not} needed.
\end{pick}
\src{ca03\_noor ``Static Instruction Scheduling'' through ``Dynamic Instruction Scheduling''; ca04 ``Multiple Issue and Static Scheduling'' through ``Superscalar v.\ VLIW''}

\section{Limits of ILP (use these to argue \emph{against} very wide issue)}
\subsection{Practical limits}
\begin{itemize}
\item \textbf{1 branch in 5 instructions:} how do you keep a 5-wide machine busy?
\item You need roughly \textbf{pipeline depth $\times$ number of FUs} independent instructions (e.g.\ $5\times4=15$--$20$) to keep the machine busy.
\item Hardware: \emph{easy} --- more instruction bandwidth, duplicate FUs; \emph{hard} --- more register-file ports; \emph{harder} --- more memory ports; wide decode hurts clock rate/pipeline depth.
\item Practical window sizes 32--126; larger windows are impractical (storage, comparisons, issue rate). Real processors are limited by number of FUs, buses, register ports.
\item \textbf{Doubling issue from 3--6 to 6--12} would need 3--4 data memory accesses/cycle, 2--3 branches resolved/cycle, renaming/accessing $>20$ registers/cycle, fetching 12--24 instructions/cycle $\Rightarrow$ the \textbf{clock rate suffers} (the widest-issue processor, Itanium 2, has the slowest clock and the highest power).
\item \textbf{Energy:} transistors switching $\propto$ peak issue rate, performance $\propto$ sustained rate; the gap grows, so \emph{multiple-issue techniques are energy-inefficient}.
\end{itemize}

\subsection{The ideal-processor study}
Assumptions: infinite renaming registers (no WAR/WAW), perfect branch and jump prediction, perfect caches (all loads/stores 1 cycle), perfect memory-address alias analysis, unbounded window and issue. ILP is then limited \emph{only by true data dependences}. Real model comparison: IBM Power5 issues 4/clk, 48 int + 40 FP rename registers, 2--6\% misprediction, 64\,KB I / 32\,KB D / 1.92\,MB L2 / 36\,MB L3.

\subsection{Head-to-head of four real ILP designs}
\begin{center}
\scriptsize
\begin{tabular}{L{2.5cm}L{3.7cm}L{1.9cm}L{1.9cm}L{1.1cm}L{1.1cm}L{1.0cm}L{1.0cm}}
\toprule
\tblhead{\textbf{Processor} & \textbf{Microarchitecture} & \textbf{Fetch/ issue/ exec} & \textbf{Int/FP units} & \textbf{GHz} & \textbf{Trans.} & \textbf{Die} & \textbf{Power}}\\
\midrule
Pentium 4 Extreme & Speculative dynamic; deeply pipelined; SMT & 3/3/4 & 7 int, 1 FP & 3.8 & 125 M & 122 mm$^2$ & 115 W \\
AMD Athlon 64 FX-57 & Speculative dynamic & 3/3/4 & 6 int, 3 FP & 2.8 & 114 M & 115 mm$^2$ & 104 W \\
IBM Power5 (1 CPU) & Speculative dynamic; SMT; 2 cores/chip & 8/4/8 & 6 int, 2 FP & 1.9 & 200 M & 300 mm$^2$ (est.) & 80 W (est.) \\
Itanium 2 & Statically scheduled, VLIW-style & 6/5/11 & 9 int, 2 FP & 1.6 & 592 M & 423 mm$^2$ & 130 W \\
\bottomrule
\end{tabular}
\end{center}
\normalsize

\begin{center}
\small
\begin{tabular}{L{3.6cm}C{2.3cm}C{2.3cm}C{2.3cm}C{2.3cm}}
\toprule
\tblhead{\textbf{Efficiency rank (1=best)} & \textbf{Itanium 2} & \textbf{Pentium 4} & \textbf{Athlon 64} & \textbf{Power5}}\\
\midrule
SPECInt / transistor & 4 & 2 & 1 & 3 \\
SPECFP / transistor & 4 & 2 & 1 & 3 \\
SPECInt / area & 4 & 2 & 1 & 3 \\
SPECFP / area & 4 & 2 & 1 & 3 \\
SPECInt / Watt & 4 & 3 & 1 & 2 \\
SPECFP / Watt & 2 & 4 & 3 & 1 \\
\bottomrule
\end{tabular}
\end{center}
\textbf{Take-aways (slide):} no single leader; Athlon leads SPECInt; Itanium 2 and Power5 dominate on SPECFP; Itanium 2 is the least efficient on everything except SPECFP/Watt; Athlon and Pentium 4 use transistors and area well; Power5 is the best energy user on SPECFP and tied on SPECInt. Itanium (VLIW/EPIC) was ``not a significant breakthrough in scaling ILP''.

\begin{dec}[How wide should the core be?]
Quote the slide conclusion: \textbf{3--6 issue is the practical limit}. Choose a core width of about 4 for a general-purpose design; go narrower for energy/area-limited designs; beyond that, spend transistors on \emph{TLP} (multithreading, more cores). Right ILP/TLP balance is ``decided in the marketplace'': servers favour TLP, desktops may still prize single-thread speed.
\end{dec}
\src{ca04 ``Limits to Multi-Issue Machines'', ``Studies of the Limitations of ILP''; ca06 (whole deck)}

\begin{trap}[ILP traps]
\begin{itemize}
\item True dependences cannot be renamed away; only WAR/WAW can.
\item Scoreboard stalls on WAR/WAW; Tomasulo does not (renaming).
\item Tomasulo alone gives \emph{imprecise} interrupts; you need the ROB for precise ones/speculation.
\item VLIW saves hardware but costs code size and binary compatibility, and one stall freezes the whole bundle.
\item Unrolling is illegal/ineffective with loop-carried dependences or unknown aliasing.
\end{itemize}
\end{trap}

\section{Chapter 3 cheat-sheet}
\begin{key}
\begin{itemize}
\item Unrolled-and-scheduled loop: 14 cycles for 4 iterations (3.5/iter) vs.\ 9 originally.
\item Tomasulo: RS + CDB + renaming; window $\le14$ vs.\ scoreboard $\le5$; Tomasulo last write cycle 57 vs.\ 62.
\item ROB: out-of-order execute, in-order commit $\Rightarrow$ precise exceptions + speculation recovery.
\item VLIW: simple hardware, big code, no binary compatibility, lock-step. Superscalar: opposite.
\item Practical issue width 3--6; need depth$\times$FUs independent instructions; wider $\Rightarrow$ slower clock, more power.
\end{itemize}
\end{key}

\chapter{Memory Hierarchy: Caches, Main Memory, Virtual Memory}\label{ch:mem}

\section{The framework: AMAT and the four levers}
\[
\AMAT = t_{hit}+\text{miss rate}\times\text{miss penalty},\qquad
\CPI_{total}=\CPI_{base}+\frac{\text{mem.\ refs}}{\text{instr.}}\times\text{miss rate}\times\text{miss penalty}\bs
\]
(The slides use exactly this CPI form in the memory-width example: $\CPI=2+1.2\times3\%\times64=4.3$.)

\begin{center}
\small
\begin{tabular}{L{3.8cm}L{11.6cm}}
\toprule
\tblhead{\textbf{Lever} & \textbf{Techniques (slides)}}\\
\midrule
Reduce \textbf{miss penalty} & Multilevel caches (the most important), critical word first \& early restart, read-miss priority over writes, merging write buffer, victim caches \\
Reduce \textbf{miss rate} & Larger blocks, larger caches, higher associativity, way prediction, pseudo-associative caches, compiler optimisations \\
Reduce penalty/rate via \textbf{parallelism} & Non-blocking caches, hardware prefetch, compiler-controlled prefetch \\
Reduce \textbf{hit time} & Small \& simple caches, avoid address translation in indexing, pipelined cache access, trace caches \\
\bottomrule
\end{tabular}
\end{center}
\textbf{Why hit time matters most:} it affects the clock cycle, so a fast hit helps \emph{everything} and its importance is ``multiplied'' beyond the AMAT formula; also miss penalty (in cycles) depends on clock cycle time.

\section{Step 1: diagnose the miss type (the ``3 C's'' + coherence)}
\small
\begin{longtable}{L{2.6cm}L{4.4cm}L{4.4cm}L{4.0cm}}
\toprule
\tblhead{\textbf{Miss type} & \textbf{Definition} & \textbf{Reduced by} & \textbf{Does NOT fix}}\\
\midrule
\endhead
Compulsory (cold) & First access to a block; measured by an infinite cache; independent of cache size & Larger blocks (spatial locality), prefetching & Bigger cache or more associativity \\
Capacity & Cache cannot hold all blocks needed; measured by fully associative cache of same size; falls as cache grows & Larger cache; blocking/tiling (shrinks working set); better temporal locality & Associativity \\
Conflict (collision) & Too many blocks map to the same set; vanishes with full associativity; falls as associativity rises & Higher associativity, victim cache, pseudo-associative, way-predicted set-assoc., code/data layout & Compulsory/capacity \\
Coherence (multiprocessor) & Block invalidated by another processor's write. \textbf{True sharing:} the same word is written by another processor. \textbf{False sharing:} different words that happen to be in the same cache line & Data layout (separate independently-updated data into different lines)\bs; smaller blocks\bs & Bigger caches \\
\bottomrule
\end{longtable}
\normalsize
\textbf{Slide examples:} 32\,B lines, array of 32 doubles: first touch of \texttt{array[0]} is a cold miss; prefetching \texttt{array[1..3]} turns later accesses into hits. In a direct-mapped cache, \texttt{array[0]}$\to$\texttt{array[16]}$\to$\texttt{array[0]} conflicts even though the cache is not full. In a multiprocessor, P2 writing \texttt{array[1]} invalidates P1's line holding \texttt{array[0]}: false-sharing miss.
\src{ca05 ``Classify Cache Misses --- 3 C's''; Cache-Misses (``Taxonomy of Cache Misses'')}

\section{Reducing miss penalty}
\subsection{Multilevel caches (``probably the best miss-penalty reduction'')}
\[
\AMAT=t_{hit,L1}+m_{L1}\big(t_{hit,L2}+m_{L2}^{local}\times P_{L2}\big),\quad
\text{global miss rate}=m_{L1}\times m_{L2}^{local}
\]
\begin{itemize}
\item \textbf{Local} miss rate $=$ misses in this cache / accesses \emph{to this cache}. \textbf{Global} $=$ misses / accesses generated \emph{by the CPU}. \textbf{Global miss rate is what matters} when judging L2/L3.
\item L1 capacity \emph{and} conflict misses are served by L2, avoiding main memory. A big L2 ($\gg$L1) gives a global miss rate close to a single-level cache of that size; the slide's conclusion is that L2 \emph{cache size} is what matters.
\item \textbf{L2 associativity:} direct-mapped hit 10\,cc, local miss 25\%; 2-way 10.1$\to$10 or 11\,cc, local miss 20\%; $P_{L2}=100$\,cc: DM $=10+0.25\times100=35$; 2-way $=10+0.20\times100=30$ (31 if the hit costs 11). \textbf{Higher associativity pays off at L2} because the extra hit time is tiny compared to the miss penalty.
\end{itemize}
\begin{ex}
1000 refs, 40 L1 misses, 20 L2 misses. $m_{L1}=4\%$; local $m_{L2}=20/40=50\%$; global $=2\%$. With $t_{L1}=1$, $t_{L2}=10$, $P_{L2}=100$, 1.5 refs/instr: $\AMAT=1+0.04(10+0.5\times100)=3.4$\,cc; stall cycles/instr $=(3.4-1)\times1.5=3.6$.
\end{ex}
\begin{pick}
Almost always. Essential when the processor is fast and memory is far (servers, desktops: L1+L2+L3). PMDs use two levels (64\,KB L1 + 256\,KB L2) to save energy/area. Hardware complexity 2 (costly if block sizes differ between levels).
\end{pick}

\subsection{Critical word first and early restart}
Do not wait for the whole block. \textbf{Critical word first:} fetch the missed word first and send it to the CPU, then fill the rest of the block. \textbf{Early restart:} fetch in normal order but restart the CPU the moment the requested word arrives. \textbf{Value depends on block size and spatial locality:} useful mainly for \emph{large blocks}; if the next access is the next sequential word in the block, the benefit shrinks.

\subsection{Read-miss priority over writes; merging write buffer}
\begin{itemize}
\item \textbf{Write-through with a write buffer:} a read miss may need data still sitting in the buffer (RAW through memory). Rather than waiting for the buffer to empty (increased read-miss penalty; one old MIPS 1000 design with a 4-word buffer saw about 50\% worse), \textbf{check the buffer for an address match and let the read go ahead if there is none}.
\item \textbf{Write-back:} on a read miss that replaces a dirty block, copy the dirty block to a buffer, do the read first, then write the dirty block: the CPU restarts sooner.
\item \textbf{Merging write buffer:} an entry holds several words but a write is usually one word. Merge a new write into an existing valid entry with a matching address: fewer full-buffer stalls and faster multiword writes. Mainly used with write-through caches.
\end{itemize}

\subsection{Victim cache}
A small \textbf{fully-associative} cache between the cache and its refill path holding only blocks \emph{just evicted}. On a miss, check it before going to the next level; if present, \textbf{swap} victim and cache block. No penalty increase because both are probed together.
\begin{itemize}
\item Jouppi: removes \textbf{20--95\% of the misses} of a \textbf{4\,KB direct-mapped} cache with 1--5 victim entries. AMD Athlon has 8 entries.
\end{itemize}
\begin{pick}
A direct-mapped (small, fast, simple) L1 suffering \emph{conflict} misses --- you keep the hit time of direct-mapped and recover much of 2-way's miss rate. Hardware complexity 2.
\end{pick}

\section{Reducing miss rate}
\subsection{Larger block size}
Reduces compulsory misses (spatial locality) but \textbf{raises miss penalty} and, in small caches, conflict and capacity misses. Pick the block that minimises \textbf{AMAT}, not miss rate.

\begin{center}
\small
\begin{tabular}{C{1.9cm}C{2.0cm}C{1.9cm}C{1.9cm}C{1.9cm}C{1.9cm}}
\toprule
\tblhead{\textbf{Block (B)} & \textbf{Penalty (cc)} & \textbf{4K} & \textbf{16K} & \textbf{64K} & \textbf{256K}}\\
\midrule
16 & 82 & 8.027 & 4.231 & 2.673 & 1.894 \\
32 & 84 & \textbf{7.082} & 3.411 & 2.134 & 1.588 \\
64 & 88 & 7.160 & \textbf{3.323} & \textbf{1.933} & \textbf{1.449} \\
128 & 96 & 8.469 & 3.659 & 1.979 & 1.470 \\
256 & 112 & 11.651 & 4.685 & 2.288 & 1.549 \\
\bottomrule
\end{tabular}
\end{center}
AMAT (cc) vs.\ block size; memory model: 80\,cc overhead + 2\,cc per 16\,B. \textbf{32--64\,B blocks dominate}; a 256\,B block only wins once the cache is 256\,KB (miss rate vs.\ block table: 4\,KB cache 256\,B block 9.51\% vs.\ 32\,B 7.24\%). Pentium 4's L2 uses 128\,B blocks.

\begin{dec}[Block size]
Small L1 $\rightarrow$ 32--64\,B. Large L2/L3 $\rightarrow$ 64--128\,B. Higher memory latency or bandwidth (overhead dominates the penalty) $\rightarrow$ somewhat larger blocks\bs. Small cache with 128/256\,B blocks $\rightarrow$ wrong (conflict and capacity misses rise).
\end{dec}

\subsection{Larger caches and higher associativity}
\begin{itemize}
\item \textbf{Larger cache:} helps capacity and conflict misses; costs hit time, area, power; popular for off-chip/L2.
\item \textbf{Associativity:} 8-way is nearly as good as fully associative. \textbf{2:1 rule:} 2-way of size $N/2$ $\approx$ direct-mapped of size $N$ (for caches $<$128\,KB). Cost: longer hit time $\Rightarrow$ clock cycle $\times1.10$ (2-way), $\times1.12$ (4-way), $\times1.14$ (8-way) relative to 1-way.
\end{itemize}
\begin{center}
\small
\begin{tabular}{C{2.4cm}C{2.2cm}C{2.2cm}C{2.2cm}C{2.2cm}}
\toprule
\tblhead{\textbf{Size (KiB)} & \textbf{1-way} & \textbf{2-way} & \textbf{4-way} & \textbf{8-way}}\\
\midrule
4 & 3.44 & 3.25 & \textbf{3.22} & 3.28 \\
8 & 2.69 & 2.58 & \textbf{2.55} & 2.62 \\
16 & \textbf{2.23} & 2.40 & 2.46 & 2.53 \\
32 & \textbf{2.06} & 2.30 & 2.37 & 2.45 \\
64 & \textbf{1.92} & 2.14 & 2.18 & 2.25 \\
128 & \textbf{1.52} & 1.84 & 1.92 & 2.00 \\
256 & \textbf{1.32} & 1.66 & 1.74 & 1.82 \\
512 & \textbf{1.20} & 1.55 & 1.59 & 1.66 \\
\bottomrule
\end{tabular}
\end{center}
AMAT with the clock-stretch penalty included (the bold cell is the best per row; in the slide the \emph{worse-than-left} cells are highlighted). \textbf{Reading:} for tiny (4--8\,KB) caches 2--4 way wins; \textbf{from 16\,KB upward direct-mapped has the lowest AMAT} once the slower clock is charged. This is the quantitative reason an L1 that sets the clock is small and low-associativity, while an L2 (off the critical path) can be highly associative.

\subsection{Way prediction}
Keep extra bits per set predicting which way will hit next; the multiplexer is set early and \textbf{only one tag comparison is done} first. A wrong prediction costs extra cycle(s) to check the remaining ways. Alpha 21264 uses it in its 2-way I-cache with $>85\%$ accuracy (SPEC95). Predictor styles: last-used way, MRU, history table. Metrics: accuracy, misprediction rate, energy reduction.
\begin{pick}
Set-associative caches where \textbf{energy and hit time} matter: PMDs, instruction caches (sequential access makes prediction accurate), cores where a full parallel tag compare is too slow/hot. It keeps most of the associativity benefit while activating one way. Cost: predictor hardware, an extra access on a misprediction.
\end{pick}

\subsection{Pseudo-associative (column-associative) cache}
Start as a \textbf{direct-mapped} cache; on a miss, probe a second location (invert the high-order index bit) \textbf{sequentially}, not in parallel. Gives fast hit (primary) and slow hit (alternate); swap blocks so the frequent one sits in the fast slot. \textbf{Problem:} a variable 1-or-2-cycle hit is hard for a CPU pipeline $\Rightarrow$ better for \textbf{caches not tied directly to the processor (L2)}; used in the L2 of the MIPS R10000 (and similar in UltraSPARC).
\begin{ex}[AMAT, slide values]
2\,KB: 1-way 5.90, 2-way 4.90, pseudo 4.95. 128\,KB: 1-way 1.50, 2-way 1.45, pseudo 1.371. Pseudo-associativity gets (nearly) 2-way miss rate at direct-mapped hit speed.
\end{ex}

\begin{center}
\small
\begin{tabular}{L{3.4cm}L{3.7cm}L{3.7cm}L{4.2cm}}
\toprule
\tblhead{\textbf{} & \textbf{Direct-mapped} & \textbf{Pseudo-associative} & \textbf{2-way set-associative}}\\
\midrule
Locations per block & 1 & 2 (sequential) & 2 (parallel) \\
Normal hit latency & Very low & Very low (primary) & Low \\
Conflict misses & High & Lower & Lower \\
Hardware & Low & Moderate & Higher \\
Energy & Low & Low/moderate & Higher \\
\bottomrule
\end{tabular}
\end{center}

\subsection{Compiler optimisations (zero hardware cost)}
\begin{itemize}
\item \textbf{Code:} reorder procedures (profile-guided) to cut conflicts; McFarling: $-50\%$ misses for a 2\,KB direct-mapped I-cache with 4\,B blocks, $-75\%$ for 8\,KB; optimised programs on a direct-mapped cache miss less than unoptimised on 8-way of the same size.
\item \textbf{Merging arrays:} combine parallel arrays (\texttt{val[]}, \texttt{key[]}) into an array of structs $\Rightarrow$ better spatial locality, fewer conflicts.
\item \textbf{Loop interchange:} change loop order so the inner loop walks memory sequentially (row-major). $1024\times1024$ ints, 64\,B lines (16 ints): \textbf{3{,}145{,}728 $\rightarrow$ 196{,}608 misses for $C=A+B$ ($-93.75\%$)}.
\item \textbf{Loop fusion:} merge loops over the same data so a value is reused while in cache (temporal locality). Example: A, B, C larger than the cache: $262{,}144 \rightarrow 196{,}608$ line transfers ($-25\%$).
\item \textbf{Blocking (tiling):} operate on $B\times B$ sub-matrices that fit in cache. Worst-case capacity misses $2N^3+N^2 \rightarrow 2N^3/B+N^2$; also helps register allocation. Slide example: $N=1024$, 32\,KB cache, $B=32$ (three tiles $=3\times4$\,KB$=12$\,KB fits): $\approx1.07$\,billion misses $\rightarrow\approx2.1$ million (simplified model).
\end{itemize}
\begin{pick}
Array/loop-dominated code with large data sets (scientific, image, matrix); embedded systems where adding hardware is impossible. \emph{Matrix multiply with arrays accessed both by rows and columns: interchange does not help $\Rightarrow$ use blocking.}
\end{pick}
\begin{avoid}
Pointer-chasing, irregular data structures (no static pattern); loops where interchange/fusion is illegal because of dependences.
\end{avoid}
\src{ca05 ``Reducing Miss Rate'' through ``Summary of Compiler Optimizations'', ``Way Prediction'', ``Pseudo-Associative Caches''; CachePerformance2 (way prediction, pseudo-associative, loop interchange/fusion/blocking numerics)}

\section{Reducing penalty/rate through parallelism}
\subsection{Non-blocking (lockup-free) cache}
Cache continues to supply hits during a miss. \textbf{Hit-under-miss} (one outstanding miss), \textbf{hit-under-multiple-miss / miss-under-miss} (overlap several misses; Pentium Pro allows 4). \textbf{Requires an out-of-order CPU} (scoreboard/Tomasulo) and \textbf{multiple memory banks}; greatly increases the cache controller complexity; must detect repeat misses to a block already being fetched.
\begin{itemize}
\item OoO CPUs can hide L1 misses that hit in L2, but \textbf{not a significant fraction of L2 misses}.
\item Example (8\,KB, FP): miss rate$\times$penalty: DM $11.4\%\times16=1.84$; 2-way $10.7\%\times16=1.71$. For integer: DM $7.4\%\times16=1.18$, 2-way $6.0\%\times16=0.96$. Hit-under-miss does \textbf{not increase hit time}, unlike associativity.
\item Hard to evaluate: effective miss penalty is the \emph{non-overlapped} time the CPU stalls.
\end{itemize}
\begin{pick}
Any out-of-order superscalar, especially memory-intensive FP code; pair with banked/interleaved memory. Useless on a strict in-order blocking CPU.
\end{pick}

\subsection{Hardware prefetching (stream buffers)}
Fetch what you expect to need before it is requested, outside the cache. Alpha AXP 21064: on an I-miss fetch 2 blocks (target to I-cache, next to the \textbf{instruction stream buffer}); 1/4/16-block ISB catches 15--25\%/50\%/72\% of misses. Data: one data stream buffer removed 25\% of a 4\,KB cache's misses, 4 streams 43\%; 8 streams 50--70\% (for two 64\,KB 4-way caches, scientific programs). \textbf{Needs spare memory bandwidth} (otherwise it hurts).
\begin{ex}[Slide]
Hit time 2\,cc, prefetch-hit time 1\,cc, prefetch-hit rate 25\%, miss rate 1.10\% (8\,KB), miss penalty 50\,cc:\\
$\AMAT=2+0.011\times0.25\times1+0.011\times0.75\times50=2.41525$\,cc. A cache without prefetching needs a miss rate of only 0.83\% (between 8\,KB at 1.10\% and 16\,KB at 0.64\%) to match.
\end{ex}

\subsection{Compiler-controlled prefetching}
Special prefetch instructions (register prefetch: HP PA-RISC; cache prefetch: MIPS IV, PowerPC, SPARC v9), non-faulting (a form of speculative execution), best for loops (e.g.\ \texttt{prefetch(b[j+7][0])}). Cost: issuing prefetches takes instruction slots, so benefit must exceed issue cost; needs a non-blocking cache. Works for instructions too.

\begin{dec}[Hardware vs.\ compiler prefetch]
Hardware (stream buffers): sequential/streaming access, no compiler support, spare bandwidth. Compiler: regular strided loops where you know the distance; embedded systems with controllable code. Skip prefetching on bandwidth-starved or energy-limited designs.
\end{dec}
\src{ca05 ``Reducing Cache Miss Penalty or Miss Rate Via Parallelism'' through ``Compiler-Controlled Prefetching''}

\section{Reducing hit time}
\begin{itemize}
\item \textbf{Small and simple caches:} the slow part of a hit is index $\rightarrow$ read tag $\rightarrow$ compare. Keep L1 small enough to be on-chip; direct-mapped is simpler/faster (more misses but faster hit); for L2 keep tags on-chip and data off-chip. Used when the cache sets the clock cycle.
\item \textbf{Avoid address translation during indexing.} \emph{Physically addressed} caches translate with the TLB first; a \emph{virtually addressed} cache indexes in parallel with translation. Problems of pure virtual caches:
 \begin{itemize}
 \item \textbf{Protection} bits must be copied on miss and checked every access.
 \item \textbf{Process switch} $\Rightarrow$ same virtual address $\neq$ same physical data $\Rightarrow$ flush (huge task-switch cost and compulsory misses); fix: \textbf{PID in the tag} (PIDs save 0.6--4.3\% miss rate over purging; cost 0.3--0.5\% for uniprocess).
 \item \textbf{Synonyms/aliases:} two virtual addresses for one physical address $\Rightarrow$ two copies; fix: hardware anti-aliasing (check PA on miss, invalidate), or \textbf{page colouring} (aliases share low-order bits, e.g.\ 18 bits).
 \item \textbf{I/O} uses physical addresses; needs mapping to VA.
 \end{itemize}
 Compromise from the slides: overlap cache access with VA translation, requiring the cache index to remain invariant across translation (trivial if the cache is small); used in Alpha 21164 and UltraSPARC III.
\item \textbf{Pipelined cache access / pipelined writes:} write hits normally need tag check before data write; pipeline into two stages so the \emph{current write's tag check overlaps the previous write's data update}; looks like one write per cycle, short cycle time. Works if the CPU does not immediately depend on the written data (watch read/write ordering). Reads are unaffected (tag check in parallel with data read).
\item \textbf{Trace cache:} stores dynamic instruction \emph{traces} from branch entry to exit, avoiding wasted space in conventional blocks (a taken branch every 5--10 instructions means much of a 16-instruction/64\,B block is unused). Costs: complicated address mapping, duplicate storage of the same instructions in different traces. Used in Pentium 4 (NetBurst).
\end{itemize}
\begin{pick}
Small simple L1 / pipelined writes: every high-clock design. Virtual/overlapped translation: tight L1 designs with a small cache. Trace cache: wide-issue front ends that need fetch bandwidth across taken branches (not energy-constrained designs).
\end{pick}
\src{ca05 ``Reducing Hit Time'' through ``Trace Caches (Cont.)''}

\section{Cache optimisation summary table (from the slides)}
\small
\begin{longtable}{L{4.6cm}C{1.2cm}C{1.2cm}C{1.2cm}C{1.8cm}L{5.2cm}}
\toprule
\tblhead{\textbf{Technique} & \textbf{Miss pen.} & \textbf{Miss rate} & \textbf{Hit time} & \textbf{HW compl.} & \textbf{Comment}}\\
\midrule
\endhead
Multilevel caches & $+$ & & & 2 & Costly; widely used \\
Critical word first / early restart & $+$ & & & 2 & Widely used \\
Read priority over writes & $+$ & & & 1 & Trivial for uniprocessor \\
Merging write buffer & $+$ & & & 1 & Used with write-through \\
Victim caches & $+$ & $+$ & & 2 & AMD Athlon: 8 entries \\
Larger block size & $-$ & $+$ & & 0 & Trivial; Pentium 4 L2 128\,B \\
Larger cache size & & $+$ & $-$ & 1 & Widely used, esp.\ L2 \\
Higher associativity & & $+$ & $-$ & 1 & Widely used \\
Way-predicting caches & & & $+$ & 2 & UltraSPARC III I-cache \\
Pseudo-associative & & $+$ & & 2 & MIPS R10000 L2 \\
Compiler techniques (miss rate) & & $+$ & & 0 & Software is the challenge \\
Non-blocking caches & $+$ & & & 3 & All out-of-order CPUs \\
HW prefetch (instr./data) & $+$ & $+$ & & 2/3 & UltraSPARC III \\
Compiler-controlled prefetch & $+$ & $+$ & & 3 & Needs non-blocking cache \\
Small and simple caches & & & $+$ ($-$ miss rate) & 0 & Trivial; widely used \\
Avoid translation in indexing & & & $+$ & 2 & Trivial if small; Alpha 21164 \\
Pipelined cache access & & & $+$ & 1 & Widely used \\
Trace cache & & & $+$ & 3 & Pentium 4 \\
\bottomrule
\end{longtable}
\normalsize
($+$ improves, $-$ hurts; HW complexity 0 = none, 3 = highest.) Note the \emph{conflicts of interest}: larger blocks hurt penalty; larger size/associativity hurt hit time; small/simple caches hurt miss rate.

\section{Main memory}
\subsection{DRAM vs.\ SRAM and key metrics}
\begin{itemize}
\item \textbf{Main memory is DRAM:} dynamic, needs periodic refresh (about 8\,ms, $\sim$1\% of time); address split into RAS and CAS halves (2-D array). \textbf{Caches are SRAM:} no refresh; 6 transistors/bit vs.\ 1. DRAM/SRAM size ratio 4--8; SRAM/DRAM cost and cycle-time ratio 8--16.
\item \textbf{Latency:} access time (request $\rightarrow$ word arrives) and cycle time (minimum between requests, $>$ access time) $\Rightarrow$ miss penalty. \textbf{Bandwidth:} bytes per unit time $\Rightarrow$ time to transfer the block (matters for I/O and large-block L2 misses).
\item DDR trend (slide tables): total best-case access time improved only from 42\,ns (DDR1, 2000) to 26\,ns (DDR4, 2016), whereas DIMM bandwidth rose from 2128 (PC2100) to 21{,}300 MiB/s (PC21300) --- again \textbf{bandwidth far outgrows latency}.
\end{itemize}

\subsection{Three ways to raise memory bandwidth}
\begin{center}
\small
\begin{tabular}{L{2.4cm}L{4.6cm}L{4.2cm}L{4.4cm}}
\toprule
\tblhead{\textbf{Organisation} & \textbf{What it is} & \textbf{Pros} & \textbf{Cons}}\\
\midrule
Simple & CPU, cache, bus, memory all one word wide & Cheapest & Block transfer is serial: penalty $\propto$ words per block \\
Wide & Memory/bus (and cache side) $N$ words wide (Alpha 64/256\,bit, UltraSPARC 512\,bit) & Bandwidth $\times N$; penalty drops accordingly & Wider bus cost; MUX between cache and CPU can be on critical path (put MUX between L1 and L2); ECC harder (partial update needs read of the rest); expansion in larger increments \\
Interleaved & Same width bus, $N$ memory banks; bank $=$ address mod $\#$banks & Overlaps bank accesses, shares bus; great for \textbf{sequential} access (cache blocks); cheap & Needs $\ge$ (access time in cc) banks to deliver one word/cycle; many small chips costly; hard to expand \\
Independent banks & Multiple controllers & Multiprocessor/I/O, non-blocking caches (hit under $n$ misses) & Controller complexity \\
\bottomrule
\end{tabular}
\end{center}

\begin{ex}[Slide: what interleaving/width buys]
Base CPI 2, 1.2 mem.\ refs/instr, miss penalty = 4 (address) + 56 (access) + 4 per word. Miss rates: 3\%/2\%/1.2\% for 1/2/4-word blocks.
\begin{center}
\small
\begin{tabular}{L{6.2cm}C{2.3cm}C{2.3cm}C{2.3cm}}
\toprule
\tblhead{\textbf{Configuration} & \textbf{1-word} & \textbf{2-word} & \textbf{4-word}}\\
\midrule
64-bit bus \& memory, no interleaving & 4.30 & 5.07 & 5.69 \\
64-bit bus \& memory, interleaved & --- & 3.63 & \textbf{3.09} \\
128-bit bus \& memory, no interleaving & --- & \textbf{3.54} & 3.84 \\
\bottomrule
\end{tabular}
\end{center}
\textit{Note:} the slide labels the narrow bus ``64-bit'' and the wide bus ``128-bit''; its arithmetic treats narrow as one word per transfer and wide as two words per transfer. \textbf{Reading:} wide/interleaved beat simple by a large margin; with larger blocks \textbf{interleaving} wins (it overlaps accesses on a narrow bus), with 2-word blocks a \textbf{128-bit} bus is marginally better. If the bus cannot be widened (pins/cost), choose interleaving.
\end{ex}
\src{ca05 ``Main Memory'' through ``Independent Memory Banks''}

\subsection{Memory hierarchy by device class (slide figures)}
\small
\begin{tabularx}{\textwidth}{L{2.3cm}L{2.9cm}L{2.8cm}L{2.8cm}X}
\toprule
\tblhead{\textbf{Class} & \textbf{L1} & \textbf{L2 / L3} & \textbf{Main memory} & \textbf{Storage}}\\
\midrule
PMD & 64\,KB, 1\,ns & L2 256\,KB, 5--10\,ns (no L3) & 1--2\,GB, 50--100\,ns & Flash 4--64\,GB, 25--50\,$\mu$s \\
Laptop & 64\,KB, 1\,ns & L2 256\,KB 3--10\,ns; L3 4--8\,MB 10--20\,ns & 4--16\,GB, 50--100\,ns & Flash 256\,GB--1\,TB \\
Desktop & 64\,KB, 1\,ns & L2 256\,KB; L3 8--32\,MB & 8--64\,GB & Flash 256\,GB--2\,TB \\
Server & 64\,KB, 1\,ns & L2 256\,KB 3--10\,ns; L3 16--64\,MB 10--20\,ns & 32--256\,GB, 50--100\,ns & Disk 16--64\,TB (5--10\,ms), Flash 1--16\,TB (100--200\,$\mu$s) \\
\bottomrule
\end{tabularx}
\normalsize
Register files: $\sim$1000\,B/300\,ps (PMD, laptop), 2000\,B (desktop), 4000\,B/200\,ps (server). \textbf{Use this to size your hierarchy:} more levels and bigger last-level cache as you move from PMD to server.

\section{Virtual memory}
\begin{itemize}
\item \textbf{Why:} many processes each with their own address space share a smaller physical memory; need \emph{protection} (a process only touches its own blocks); faster program start (not everything resident); frees the programmer from overlays.
\item \textbf{Pages vs.\ segments:} pages $=$ fixed-size blocks (4--8\,KB), one address (page number $+$ offset), \textbf{no external fragmentation}, allocation anywhere (non-contiguous), no compaction; segments $=$ variable size (1\,B up to $2^{16}$--$2^{32}$\,B), address needs 2 words (segment number + offset). \textbf{Paging is the default choice}.
\item \textbf{Translation:} virtual page number indexes the \textbf{page table} (valid bit, physical page number); valid $=0\Rightarrow$ \textbf{page fault} (OS loads from disk). The page table is in memory, so every access would cost two memory accesses $\Rightarrow$ \textbf{TLB}, a special cache of translations.
\item \textbf{TLB:} usually small ($\le128$--256 entries); high-end: fully associative; mid-range: small $n$-way set associative. Fields: valid, dirty, tag, physical page number. Physical address then drives cache tag compare (TLB $\rightarrow$ cache interaction).
\item Example split: 32-bit VA $=$ 20-bit VPN $+$ 12-bit offset (4\,KB pages).
\end{itemize}
\src{ca05 ``Virtual Memory'' through ``Techniques for Fast Address Translation''}

\section{Scenario $\rightarrow$ memory design, quick reference}
\small
\begin{longtable}{L{5.2cm}L{4.1cm}L{6.3cm}}
\toprule
\tblhead{\textbf{Symptom / constraint} & \textbf{Choose} & \textbf{Reason}}\\
\midrule
\endhead
Small direct-mapped L1, ping-pong conflicts & Victim cache; or pseudo-assoc./2-way; or array merging/padding & Conflict misses; victim cache keeps DM hit time. \\
Streaming/sequential large arrays & Stream-buffer prefetch, moderate block (64\,B), interleaved memory, loop interchange & Compulsory/spatial misses; bandwidth is plentiful. \\
Large matrix multiply & Blocking (tiling) with 3 tiles $\le$ cache & Capacity misses; interchange cannot help both row/column accesses. \\
OoO core with many outstanding loads & Non-blocking cache + banked memory & Overlap misses; otherwise OoO core is wasted. \\
L1 hit time limits clock & Small, simple (DM) L1, pipelined writes, VA/translation overlap & Hit time is multiplied across every instruction. \\
PMD / battery & Two cache levels, way prediction, moderate sizes, no trace cache, selective prefetch & Energy and area; way prediction cuts tag-compare energy. \\
Server with large working set & L3 (16--64\,MB), high-associativity L2, interleaved/banked memory & Global miss rate matters; capacity misses dominate. \\
Write-through L1 & Merging write buffer + read priority over writes & Prevent buffer-full stalls and RAW delays. \\
Large blocks and long block fill & Critical word first \& early restart & CPU resumes when the needed word arrives. \\
Instruction-fetch bound, taken branches every 5--10 instr. & Trace cache or I-stream buffer; BTB & Fetch bandwidth across taken branches. \\
Multicore sharing data & Watch coherence (true/false sharing) misses; separate independently-written data into different lines\bs & Coherence misses are a distinct fourth C. \\
Cache must be tiny but fast & Direct-mapped, small block & Table: DM wins AMAT $\ge$16\,KB; for 4--8\,KB use 2--4 way. \\
\bottomrule
\end{longtable}
\normalsize

\begin{trap}[Memory traps]
\begin{itemize}
\item Use \textbf{global} miss rate (not local) to judge L2/L3. A ``50\% L2 miss rate'' can still be a good L2.
\item Larger blocks lower miss rate but raise miss penalty; pick the minimum \textbf{AMAT}, not miss rate (32--64\,B).
\item Higher associativity lowers miss rate but lengthens the cycle: direct-mapped wins AMAT for $\ge$16\,KB in the slide table.
\item Non-blocking caches are only valuable with an out-of-order CPU and multiple banks.
\item Hardware prefetch needs unused bandwidth; compiler prefetch costs instruction slots.
\item Pseudo-associative = sequential probing (fast/slow hit), not 2-way parallel; not good for an L1 tied to the pipeline.
\item Loop interchange helps only if inner loop currently strides across rows; blocking is for mixed row/column access.
\item Virtual caches need PID tags and anti-aliasing; I/O still uses physical addresses.
\end{itemize}
\end{trap}

\section{Chapter 4 cheat-sheet}
\begin{key}
\begin{itemize}
\item $\AMAT=t_{hit}+m\times P$; two-level: $t_1+m_1(t_2+m_2P_2)$; global $=m_1m_2$.
\item Block: 32--64\,B optimum. Associativity: DM best for $\ge$16\,KB (with 10/12/14\% clock penalty).
\item Victim cache: 20--95\% fewer misses in a 4\,KB DM cache. Way prediction $>85\%$ accurate.
\item Interchange $-93.75\%$, fusion $-25\%$, blocking: $2N^3+N^2\to2N^3/B+N^2$.
\item Interleaved memory: bank $=$ addr mod banks; need banks $\ge$ access cycles.
\item Prefetch needs spare bandwidth; non-blocking needs OoO + banks.
\end{itemize}
\end{key}

\chapter{Thread-Level Parallelism: Multithreading, SMT, Multicore}\label{ch:tlp}

\section{Decision map}
\small
\begin{tabularx}{\textwidth}{L{4.4cm}L{4.2cm}X}
\toprule
\tblhead{\textbf{Scenario} & \textbf{Feature} & \textbf{Why}}\\
\midrule
Throughput server, many independent threads/requests & Multicore (CMP) with multithreading per core & Applications with TLP (DB, web) map directly; ILP alone is saturated. \\
Pipeline stalls are \emph{long and rare} (e.g.\ L2 misses), want no slowdown of a single thread & Coarse-grained multithreading & Switch only on costly stalls; pipeline refill $\ll$ stall time. \\
Stalls are \emph{short and frequent}; hardware can switch every cycle & Fine-grained multithreading & Hides short and long stalls; interleaves each cycle. \\
Out-of-order superscalar with idle issue slots; want cheap extra throughput & SMT & Reuses renaming and OoO hardware; fills unused slots from another thread. \\
Single important thread must run as fast as possible & One big (SMT-capable) superscalar core; preferred-thread priority & Multi-core does not speed one thread; SMT with priority protects it. \\
Limited die area / thermal budget but threaded software & Several small cores & Each core simpler and easier to design/manufacture; better energy per throughput. \\
Applications hard to parallelise & Do \emph{not} rely on cores/threads & TLP gives no gain; spend on ILP/caches. \\
\bottomrule
\end{tabularx}
\normalsize

\section{Why TLP at all}
\begin{itemize}
\item ILP gives implicit parallelism within a loop or basic block; it saturates at 3--6 issue (Chapter~\ref{ch:ilp}), and for most applications \textbf{most execution units sit idle} even on an 8-way superscalar (Tullsen et al.).
\item \textbf{TLP is explicit:} multiple threads, each with its own instructions, data, PC and register state. Goals: improve \emph{throughput} of machines running many programs and \emph{execution time} of multithreaded programs. TLP ``could be more cost-effective to exploit than ILP''.
\item Natural parallelism: online transaction systems (many independent queries/updates), data mining, databases, scientific codes; server-per-client threads; games (AI, graphics, physics as threads).
\item Industry move: IBM Power4 (2000) was the first commercial single-chip general-purpose multiprocessor (2 Power3 + shared L2); Sun, AMD, Intel followed.
\end{itemize}
\src{ca06 ``Performance beyond single thread ILP'', ``Thread Level Parallelism (TLP)''; ca08 ``Why multi-core?'', ``Thread-level parallelism''}

\section{Multithreading on one core}
Multiple threads \textbf{share the functional units} of one processor by overlapping. The processor \emph{duplicates the independent state of each thread}: a separate register file, a separate PC, and (for independent programs) a separate page table. Memory is shared via the virtual-memory mechanism. Hardware thread switch is \emph{much faster} than a full process switch (100s--1000s of clocks).

\begin{center}
\small
\begin{tabular}{L{2.4cm}L{4.2cm}L{4.5cm}L{4.4cm}}
\toprule
\tblhead{\textbf{} & \textbf{Fine-grained} & \textbf{Coarse-grained} & \textbf{SMT (simultaneous)}}\\
\midrule
When it switches & Every instruction/clock, round-robin, skipping stalled threads & Only on costly stalls (e.g.\ L2 miss) & No switching: instructions from several threads issue in the \emph{same cycle} \\
Hides & Short \emph{and} long stalls & Only long stalls (pipeline refill $\ll$ stall time) & Short and long stalls \emph{and} wasted issue slots \\
Single-thread cost & Slows a ready thread (delayed by others) & Does not slow a thread (others run only when it stalls) & Can slow a thread unless priority is given \\
Hardware & Switch every cycle; per-thread state & No very fast switch needed; pipeline must be flushed/frozen on switch & OoO superscalar + per-thread rename table, PC, (logically) per-thread ROB \\
Example & Sun Niagara & IBM AS/400 & Pentium 4 (HT), IBM Power5 \\
\bottomrule
\end{tabular}
\end{center}

\begin{dec}[Pick the multithreading style]
\begin{itemize}
\item \textbf{Coarse-grained} if the core is simple/in-order, stalls are dominated by rare, long events (L2/L3 misses), and single-thread speed must not degrade; avoid when short stalls dominate (refill cost cannot be recovered).
\item \textbf{Fine-grained} for throughput-oriented cores with many threads where single-thread latency is secondary (server/web/DB); needs per-cycle switching hardware.
\item \textbf{SMT} when you already have a wide out-of-order superscalar (renaming, large physical register set, OoO completion): \emph{the marginal hardware is small}. Not worthwhile for narrow in-order cores.
\end{itemize}
\end{dec}

\subsection{Why SMT is cheap on an OoO core}
Dynamically scheduled processors already have: a large set of virtual (physical) registers that can hold several threads' register sets; \textbf{register renaming} gives unique register ids so instructions from different threads mix in the datapath without confusing sources and destinations; \textbf{out-of-order completion} lets threads proceed independently. You only add a per-thread renaming table and separate PCs; independent commit is obtained by logically keeping a separate ROB per thread.

\subsection{SMT design challenges and costs (IBM Power5 case study)}
\begin{itemize}
\item \textbf{Single-thread performance:} SMT needs fine-grained scheduling; a \textbf{preferred-thread} approach protects one thread but loses some throughput when that thread stalls. Power5 gives hardware-controllable \emph{relative thread priority}; with balanced threads both run slower than if each owned the machine.
\item Larger register file for multiple contexts; do not lengthen the clock cycle (more candidate instructions to issue, harder choice at commit); avoid cache/TLB conflicts between threads.
\item \textbf{Why only 2 threads in Power5?} With 4, one of the shared resources (physical registers, cache, memory bandwidth) would become a bottleneck.
\item \textbf{Changes made in Power5 to support SMT:} higher associativity in the L1 I-cache and instruction translation buffers; per-thread load/store queues; larger L2 (1.92 vs.\ 1.44\,MB) and L3; separate instruction prefetch and buffering per thread; physical registers 152 $\rightarrow$ 240; larger issue queues. \textbf{Power5 core is $\approx$24\% larger than Power4's} because of SMT.
\end{itemize}

\subsection{SMT performance data to quote}
\begin{center}
\small
\begin{tabular}{L{6.2cm}L{9.2cm}}
\toprule
\tblhead{\textbf{Measurement} & \textbf{Result (slides)}}\\
\midrule
Pentium 4 Extreme SMT, SPECrate (same benchmark $\times$ copies) & 1.01 (SPECint\_rate), 1.07 (SPECfp\_rate) \\
Pentium 4, 26 SPEC benchmarks paired with each other (262 runs) & Speedups 0.90--1.58, \textbf{average 1.20} \\
Power5, 8-processor server & 1.23 (SPECint\_rate), 1.16 (SPECfp\_rate) with SMT \\
Power5, two copies of each application & 0.89--1.41; most gained; \textbf{FP applications had the most cache conflicts and the least gains} \\
General statement (multicore deck) & ``Enables better threading (e.g.\ up to 30\%)'' \\
\bottomrule
\end{tabular}
\end{center}
\textbf{Reading:} SMT gives a modest, workload-dependent gain (roughly 1.0--1.4$\times$) and can \emph{hurt} (speedup $<$1) when threads fight over caches. It is \emph{not} a true parallel processor: there is a single copy of each resource, and threads cannot use the same functional unit simultaneously.

\begin{ex}[Area-efficiency argument for SMT vs.\ a second core]\bs
Power5 SMT cost $\approx$24\% extra core area for $\approx$16--23\% throughput. A second core costs about 100\% extra core area but can give up to $\sim$2$\times$ throughput for fully parallel software. So SMT is roughly \emph{area-neutral} on throughput but needs little extra design and shines when threads are stall-bound; a second core wins when software is truly parallel and area/power allow.
\end{ex}
\src{ca06 ``New Approach: Multithreaded Execution'' through ``Initial Performance of SMT''; ca08 ``Simultaneous multithreading (SMT)''}

\section{Multicore (chip multiprocessor, CMP)}
\begin{itemize}
\item Several cores on one processor socket; they run in parallel; \textbf{within each core threads are time-sliced} as on a uniprocessor. The OS sees each core as a separate processor and schedules threads/processes across them.
\item \textbf{Why multicore:} single-core clock frequencies cannot rise further; limited ILP; many applications are already multithreaded; general trend toward parallelism.
\item A multicore is a \textbf{shared-memory MIMD multiprocessor on one chip}: different cores run different instruction streams on different data, all sharing the same memory.
\item \textbf{Applications that benefit:} database servers, web servers, compilers, multimedia, scientific, CAD/CAM; ``anything that can be threaded''. Some applications are difficult to parallelise.
\end{itemize}

\subsection{Multicore vs.\ SMT, and combining them}
\begin{center}
\small
\begin{tabular}{L{3.2cm}L{6.0cm}L{6.2cm}}
\toprule
\tblhead{\textbf{} & \textbf{Multicore} & \textbf{SMT}}\\
\midrule
Resources & Each core has its own copy of resources & Single copy of each resource shared by threads \\
Core size/design & Several smaller, less powerful cores; easier to design and manufacture & One large, fast superscalar core \\
Strength & Great with thread-level parallelism; independent execution & Great single-thread performance; still mostly ILP-based \\
Weakness & Does nothing for one thread; coherence traffic & Not true parallelism; shared-resource contention; $\le$30\% gain \\
Threads & One (or more) per core & 2, 4 or sometimes 8 per core (Intel ``hyper-threads'') \\
\bottomrule
\end{tabular}
\end{center}
Four combinations: single-core non-SMT (standard uniprocessor), single-core SMT, multi-core non-SMT, \textbf{multi-core with SMT} (the best-throughput option; e.g.\ SMT dual-core runs four threads concurrently).

\subsection{Memory hierarchy in threaded designs}
\begin{itemize}
\item SMT only $\Rightarrow$ all caches shared. Multicore $\Rightarrow$ \textbf{L1 caches private}; L2 \textbf{private in some designs, shared in others}; \textbf{main memory always shared}. Example: dual-core hyper-threaded Xeon with private L1 and \textbf{shared} L2.
\item \textbf{Shared L2:} threads share capacity dynamically and share data without duplication\bs; risk of interference. \textbf{Private L2:} isolation and simple timing\bs; duplicated data, wasted capacity if load is unbalanced\bs. In either case sharing introduces \textbf{coherence misses} (true/false sharing; Section on miss types).
\end{itemize}

\begin{pick}[Multicore + SMT]
Server/throughput workloads and mixed multitasking desktops: plenty of independent threads, area available, TLP-friendly software. Add SMT to each core if the core is a wide OoO and memory stalls are common.
\end{pick}
\begin{avoid}
Single-threaded or hard-to-parallelise applications (give one fast core instead); extremely tight power budgets that cannot afford duplicate structures or shared-cache contention.
\end{avoid}

\begin{dec}[ILP vs.\ TLP balance (slide conclusion)]
There is no universal answer: \emph{server} market $\Rightarrow$ more TLP (many cores/threads); \emph{desktop} where single-thread speed is still a primary requirement $\Rightarrow$ keep a strong ILP core but add a few cores. ``Balance of ILP and TLP decided in marketplace.''
\end{dec}

\begin{trap}[TLP traps]
\begin{itemize}
\item SMT threads cannot use the \emph{same} functional unit at the same time; it is not a doubling of resources.
\item Coarse-grained multithreading is useless when stalls are short (pipeline refill cost dominates); fine-grained slows individual threads.
\item More threads $\ne$ more speed: Power5 stops at 2 threads because shared resources saturate; SMT can give speedup $<1$ on cache-conflict-heavy (FP) code.
\item Multicore does not accelerate a single thread; the OS must have threads to schedule.
\item False sharing causes coherence misses even though data is not shared.
\end{itemize}
\end{trap}
\src{ca08 (multicore deck: ``Multi-core architectures'' through ``Designs with private L2 caches''); ca06 ``Summary'', ``Conclusions''; Cache-Misses (coherence misses)}

\section{Chapter 5 cheat-sheet}
\begin{key}
\begin{itemize}
\item Fine-grained: switch every cycle (Niagara); coarse-grained: switch on long stalls (AS/400); SMT: same-cycle issue from several threads on an OoO core (P4, Power5).
\item SMT = per-thread PC + rename table (+ ROB per thread) on existing OoO hardware; Power5 core $+24\%$ area; typical gain 1.0--1.4$\times$ (avg 1.20 on P4; 1.16--1.23 on Power5).
\item Power5 limited to 2 threads (shared-resource bottleneck).
\item Multicore: private L1, L2 shared or private, memory shared; best with TLP; cores smaller and easier to build.
\end{itemize}
\end{key}

\chapter{Putting It Together: Scenarios, Dependencies, Checklist}\label{ch:scen}

\section{Which features only work together (check your design for consistency)}
Exam answers lose marks for choosing a feature whose prerequisite is missing. Use this table as a consistency check.

\small
\begin{longtable}{L{4.6cm}L{5.4cm}L{5.4cm}}
\toprule
\tblhead{\textbf{If you choose\ldots} & \textbf{You also need\ldots} & \textbf{Otherwise\ldots}}\\
\midrule
\endhead
Non-blocking cache (hit-under-miss) & Out-of-order CPU (scoreboard/Tomasulo), multiple memory banks & Blocking CPU stalls anyway; cannot support multiple outstanding misses \\
Hardware speculation & Dynamic branch prediction, reorder buffer (commit in order), recovery/flush & No way to undo wrong-path work; imprecise exceptions \\
Wide issue (superscalar) & Accurate predictor + BTB, more register-file/memory ports, renaming, enough independent instructions & Issue slots idle; clock rate and power suffer \\
SMT & Out-of-order superscalar core, renaming, larger physical register file, per-thread PC/rename/ROB, cache associativity/size increase & Needless on narrow in-order core; shared-resource conflicts \\
Predict-taken static scheme & Branch target known early (BTB/early BTA) & No benefit \\
VLIW & Strong compiler (alias analysis, unrolling, scheduling), no need for binary compatibility & Poor code size/utilisation, lock-step stalls \\
Loop unrolling (compiler) & More registers, renaming, no loop-carried dependence & Register pressure; false dependences \\
Hardware prefetch & Spare memory bandwidth & Prefetches compete with demand misses \\
Compiler prefetch & Non-blocking cache; issue slots & Wasted instructions; stalls \\
Way prediction / pseudo-associative & Set/column structure; tolerance for variable hit latency (better for L2 than L1 for pseudo) & Pipeline problems when hit time varies \\
Virtual-address cache & PID tags, anti-aliasing (page colouring), protection bits & Flushes on context switch; synonyms \\
Delayed branch & Compiler that can fill slots (single slot typical) & NOPs fill slots; deeper pipes are hard \\
Larger L2/L3 & Evaluate with \emph{global} miss rate & Mis-judging gains \\
Multicore & Software with threads; coherence handling & Idle cores \\
\bottomrule
\end{longtable}
\normalsize

\section{Reusable full-answer skeleton}
\begin{enumerate}[leftmargin=*]
\item \textbf{Class and goal:} (e.g.\ ``PMD: energy-efficient, real-time'').
\item \textbf{Workload analysis:} parallelism type, branch behaviour, memory behaviour, latency tolerance.
\item \textbf{Core:} pipeline depth, issue width, scheduling (static/dynamic/speculative), why.
\item \textbf{Control:} branch predictor type and size, BTB, return stack.
\item \textbf{Memory:} L1 (size, associativity, block, hit time), L2/L3, miss-penalty/miss-rate/hit-time techniques, main memory organisation, TLB.
\item \textbf{Parallelism beyond ILP:} threads per core, number of cores.
\item \textbf{Quantify at least one claim} using AMAT / CPI / stalls-per-instruction.
\item \textbf{Costs and rejected alternatives.}
\end{enumerate}

\section{Worked scenarios (model answers)}

\subsection*{Scenario 1 --- Battery-powered smartphone SoC}
\textbf{Constraints:} battery life, thermal limit, real-time UI/media, low cost, mixed integer code, occasional bursts.

{\small
\begin{longtable}{L{3.2cm}L{5.6cm}L{6.6cm}}
\toprule
\tblhead{\textbf{Block} & \textbf{Choice} & \textbf{Why}}\\
\midrule
\endhead
Class/goal & PMD: energy efficiency + real-time & Slide emphasis. Optimise energy per task, not clock. \\
Core & Narrow (2-way) in-order or modest OoO core, shallow pipeline (5--8 stages) & ILP techniques are energy-inefficient (transistors switch with peak issue rate, performance follows sustained rate); deep pipelines magnify branch penalty. \\
Branch & Small 2-bit BHT (+ small BTB, return stack) or BTFNT/static & Cheap and accurate on loops; avoids large predictor area/energy. \\
L1 & 32--64\,KB, 2-way with \textbf{way prediction}, 32--64\,B blocks & Way prediction does one tag comparison (energy); 32--64\,B minimise AMAT. \\
L2 & 256\,KB, no L3 & Slide PMD hierarchy: L1 64\,KB $\rightarrow$ L2 256\,KB $\rightarrow$ 1--2\,GB DRAM $\rightarrow$ flash. \\
Memory & Interleaved (or modestly wide) DRAM; low-power DRAM states & Bandwidth at low cost; DRAM low-power states listed in slides. \\
Power & DVFS, power gating, turn off idle cores & Dynamic power $\tfrac12CV^2f$; static power cut by gating. \\
Parallelism & Several small cores (big.LITTLE-like)\bs; \textbf{no SMT} & Small cores are easy to build; SMT adds area/power for $\sim$20\%. \\
\bottomrule
\end{longtable}}
\textbf{Rejected:} wide speculative core (power), trace cache (area/energy), large L3 (not in PMD hierarchy), VLIW (binary compatibility for apps).

\subsection*{Scenario 2 --- High-throughput web/database server}
\textbf{Constraints:} thousands of independent client requests, availability 24/7, memory-bound database queries, cost per request.

{\small
\begin{longtable}{L{3.2cm}L{5.6cm}L{6.6cm}}
\toprule
\tblhead{\textbf{Block} & \textbf{Choice} & \textbf{Why}}\\
\midrule
\endhead
Class/goal & Server: availability, scalability, throughput & Request-level parallelism: threads are independent. \\
Cores & \textbf{Many cores (CMP) + 2--4 SMT/fine-grained threads per core} & TLP is abundant; ILP of one thread is limited (most units idle). DB/web are listed as the best multi-core beneficiaries. \\
Core style & Moderate-width (4-issue) OoO with speculation, or simpler fine-grained cores & Cores are smaller/easier to build; throughput over latency. \\
Multithreading & Fine-grained or SMT & Hides short and long (DRAM) stalls; DB workloads are memory-bound. \\
L1/L2/L3 & Private L1; L2 private or shared; \textbf{large shared L3 (16--64\,MB)} & Slide server hierarchy; global miss rate matters; high-associativity L2 is cost-effective. \\
Memory & Large DRAM (32--256\,GB), \textbf{interleaved/independent banks}, non-blocking caches & Many outstanding misses from many threads; banks support multiprocessor traffic. \\
Dependability & ECC memory, redundancy; availability $=$ MTTF/(MTTF+MTTR) & Server emphasis on availability. \\
Coherence & Aware of true/false sharing & Cores share data; coherence misses appear. \\
\bottomrule
\end{longtable}}
\textbf{Rejected:} very wide single core (ILP saturates), coarse-grained MT only (short stalls), aggressive deep speculation burning power per request.

\subsection*{Scenario 3 --- Desktop/gaming: single-thread speed is the priority}
\textbf{Constraints:} one dominant thread, integer + some FP, branchy code, price-performance.

{\small
\begin{longtable}{L{3.2cm}L{5.6cm}L{6.6cm}}
\toprule
\tblhead{\textbf{Block} & \textbf{Choice} & \textbf{Why}}\\
\midrule
\endhead
Core & \textbf{4-wide dynamically scheduled speculative (Tomasulo + ROB)}, moderate-deep pipeline & Needs ILP and precise exceptions; 3--6 issue is the practical sweet spot (Athlon led SPECInt with 3-issue). \\
Branch & \textbf{Two-level / gshare predictor + BTB + return-address stack}, integrated fetch unit & Branchy code; deep pipe $\Rightarrow$ misprediction cost high. \\
L1 & Small (32--64\,KB), low associativity, pipelined writes, 64\,B block & Hit time sets the clock; DM/low-way wins AMAT $\ge$16\,KB. \\
L2/L3 & 256\,KB L2 (higher associativity) + 8--32\,MB L3 & Desktop hierarchy; L2 associativity cheap relative to penalty. \\
Latency hiding & Non-blocking cache (hit under multiple misses), HW prefetch & Out-of-order core can overlap L1-miss/L2-hit. \\
Parallelism & 2--4 cores; SMT optional (with preferred-thread priority) & Games/AI/physics in threads, but single thread still primary. \\
\bottomrule
\end{longtable}}
\textbf{Rejected:} VLIW (irregular integer code, compatibility), coarse-grained MT (no gain for one thread), very wide issue (clock/power).

\subsection*{Scenario 4 --- Scientific kernel engine (regular FP loops, huge arrays)}
\textbf{Constraints:} floating-point throughput, long-latency FP ops (add 3, mult 6, divide 24/25), matrices much larger than cache, predictable loops.

{\small
\begin{longtable}{L{3.2cm}L{5.6cm}L{6.6cm}}
\toprule
\tblhead{\textbf{Block} & \textbf{Choice} & \textbf{Why}}\\
\midrule
\endhead
FP units & Fully pipelined adder (4-stage) and multiplier (7-stage); \textbf{pipelined or duplicated divider} if divides are frequent & Initiation interval 1 for add/mult; divider with interval 25 would serialise divides. \\
Scheduling & Compiler loop unrolling + software pipelining on a VLIW/static design, \emph{or} Tomasulo with renaming & Independent iterations, known latencies: unrolling reaches 3.5 cycles/iter vs 9. Tomasulo is better if latencies vary (cache misses). \\
Branches & BTFNT or 2-bit BHT; delay slot from target & Loop back-edges almost always taken. \\
Cache software & Blocking, interchange, fusion, array merging & Capacity/conflict misses: interchange cut misses 93.75\%; blocking reduces capacity misses to $2N^3/B$. \\
Cache hardware & Stream-buffer prefetch, non-blocking cache, 64\,B blocks & Sequential/stride access, plentiful bandwidth. \\
Memory & Wide or heavily interleaved banks & Bandwidth-bound; bank count $\ge$ access cycles. \\
Parallelism & Multiple cores (DLP across data), SIMD/vector\bs{} & Same operation on lots of data. \\
\bottomrule
\end{longtable}}
\textbf{Rejected:} large BTB/predictor budget (loops are easy), heavy speculation machinery (little control-flow irregularity), SMT (FP code showed least SMT gain and most cache conflicts).

\subsection*{Scenario 5 --- Low-cost embedded controller}
\textbf{Constraints:} price is the emphasis, fixed firmware, modest performance, limited power.

\textbf{Design:} classic 5-stage in-order pipeline with forwarding and \textbf{early branch resolution in ID} (1-cycle penalty); \textbf{static} prediction (predict-not-taken, or BTFNT/profile-driven) plus \textbf{one delay slot} filled by the compiler (about 50\% useful); small direct-mapped I- and D-caches (small and simple $\Rightarrow$ fast hit, low area), maybe a single level; software techniques (loop unrolling, loop interchange, blocking) since the firmware is compiled offline; no SMT, no multiple issue. If a compiler team and fixed firmware exist and the application is a regular numeric kernel, \textbf{VLIW} reduces control hardware, but only because \emph{binary compatibility is irrelevant}.\\
\textbf{Why:} price emphasis $\Rightarrow$ minimal hardware; static techniques need no predictor memory; software cache optimisations cost zero hardware. \textbf{Rejected:} Tomasulo+ROB, BTB, trace cache, victim cache (area).

\subsection*{Scenario 6 --- Product family that must run legacy binaries across generations}
\textbf{Constraint:} next-generation chips will have more functional units and different latencies; customers keep the old binaries.\\
\textbf{Answer:} \textbf{dynamically scheduled superscalar}. Hardware finds parallelism at run time (RAW by forwarding, WAR/WAW by renaming), so the same binary benefits from more units. A VLIW compiled for $K$ units would need recompilation for $K{+}5$ (or an emulator that hurts performance), has no hazard detection, and stalls in lock-step; EPIC/IA-64 ``promises'' compatibility but remained larger in code size and was ``not a significant breakthrough'' (Itanium 2: widest issue, slowest clock, highest power).

\subsection*{Scenario 7 --- Memory-bound, irregular workload (graph analytics / pointer chasing)}
\textbf{Constraints:} huge working set, poor spatial/temporal locality, unpredictable addresses, long DRAM latency.

\textbf{Design:} Software cache optimisations (blocking, interchange) and stream prefetching don't apply to irregular access, so \textbf{hide latency with threads}: SMT/fine-grained multithreading (many threads per core) on an OoO core with a \textbf{non-blocking cache and many memory banks} (miss-under-miss); large last-level cache; victim cache for conflict misses. Coarse-grained MT is suitable only if stalls are dominated by rare long L2 misses with short refill cost. Note the slide caveat that an OoO CPU can hide L1-miss/L2-hit but \emph{not} most L2 misses, hence thread-level latency hiding. Remember SMT gains are modest and cache-conflict-sensitive (increase associativity/size as Power5 did).

\subsection*{Scenario 8 --- Numeric cache comparison (practice, constructed numbers)}
\textbf{Problem.} L1 miss penalty 40\,cc, hit time 1\,cc. Option A: 16\,KB direct-mapped, miss rate 3.0\%. Option B: 16\,KB 2-way, miss rate 2.6\%, clock stretched by 10\% (slide: 2-way $\times1.10$). Option C: option A plus a 4-entry victim cache removing 40\% of misses (slide range 20--95\%). Which is best?

\begin{ex}[Solution (AMAT in direct-mapped cycles)]
A: $1+0.030\times40=2.20$. B: $1.10+0.026\times40=2.14$ --- only 2.7\% better \emph{in AMAT}, but the 10\% longer clock also slows \emph{every non-memory instruction}, so B is a net loss. C: $1+0.030\times0.6\times40=1.72$ (assuming the victim probe hides inside the normal miss path, per the slide). \textbf{Choose C}: keeps the direct-mapped hit time (clock) and removes conflict misses cheaply. Quote the slide table: at 16\,KB, 1-way has the lowest AMAT.
\end{ex}

\subsection*{Scenario 9 --- Numeric branch-handling comparison (practice, constructed numbers)}
\textbf{Problem.} Deep pipeline (R4000-like): branch penalty 3 cycles, 20\% of instructions are branches, 60\% of branches are taken, base CPI 1.0. Compare (a) stall always, (b) predict-not-taken, (c) 2-bit BHT+BTB 90\% accurate, (d) correlating/two-level predictor 95\% accurate.

\begin{ex}[Solution]
\begin{center}
\small
\begin{tabular}{L{6cm}C{4.2cm}C{2.4cm}}
\toprule
\tblhead{\textbf{Scheme} & \textbf{Extra CPI} & \textbf{CPI}}\\
\midrule
(a) Stall always & $0.20\times3=0.60$ & 1.60 \\
(b) Predict not taken (wrong when taken) & $0.20\times0.60\times3=0.36$ & 1.36 \\
(c) 2-bit BHT + BTB, 90\% & $0.20\times0.10\times3=0.06$ & 1.06 \\
(d) Two-level, 95\% & $0.20\times0.05\times3=0.03$ & 1.03 \\
\bottomrule
\end{tabular}
\end{center}
Speedup of (c) over (a): $1.60/1.06\approx1.5$; (d) over (c): $\approx1.03$. \textbf{Conclusion:} in a deep pipe the dynamic predictor is mandatory; the extra complexity of (d) is justified only if area allows (diminishing returns), or if branch frequency is higher or the pipeline deeper.
\end{ex}

\subsection*{Scenario 10 --- Area-limited choice: big core, SMT, or multicore?}
\textbf{Constraint:} throughput workload, fixed die area, threads are available but half the time stalled on memory; one thread must stay responsive.

\begin{ex}[Derived reasoning]\bs
Using slide data: SMT $\approx+24\%$ core area for $\approx1.16$--$1.23\times$ throughput (Power5) $\Rightarrow$ throughput/area $\approx0.94$--$0.99$ (about neutral). A second core $\approx+100\%$ area for up to $2\times$ when software is fully parallel (throughput/area $\approx1.0$, but little gain if software is not parallel). \emph{Choice:} if the software is highly threaded and area allows $\rightarrow$ multicore; if area is tight and threads are stall-bound, or a thread needs priority $\rightarrow$ SMT with preferred-thread control; both (multicore + SMT) for servers with big area budgets. Widen a single core only to the 3--6 issue sweet spot: beyond that the energy and clock costs dominate.
\end{ex}

\section{Conflict resolution: when two constraints disagree}
\begin{center}
\small
\begin{tabular}{L{5.2cm}L{10.2cm}}
\toprule
\tblhead{\textbf{Conflict} & \textbf{Resolution}}\\
\midrule
High clock vs.\ low CPI & Deeper pipeline raises clock but adds load/branch delay; compensate with prediction and forwarding. Decide by the \emph{product} IC$\times$CPI$\times T_{clk}$. \\
Low miss rate vs.\ fast hit & Large/associative cache cuts misses but stretches the clock; use small DM L1 + big associative L2; victim cache/way prediction as compromises. \\
Throughput vs.\ single-thread latency & Fine-grained MT/SMT raise throughput but slow a thread; add priority or fewer threads; or heterogeneous cores. \\
Performance vs.\ power & ILP/speculation/prefetch cost energy; prefer TLP/multicore and DVFS; energy per task is the right metric. \\
Compiler simplicity vs.\ hardware simplicity & Dynamic scheduling = simple compiler, complex hardware; VLIW = opposite. Choose by compatibility and cost. \\
Bandwidth vs.\ latency & Buy bandwidth (wide/interleaved memory, banks); hide latency (OoO, non-blocking, MT, prefetch). \\
Compatibility vs.\ simplicity & Compatibility $\Rightarrow$ superscalar; fixed product $\Rightarrow$ VLIW. \\
\bottomrule
\end{tabular}
\end{center}

\section{Final ten-minute exam checklist}
\begin{enumerate}[leftmargin=*]
\item Did I state the machine class and its priority (energy / cost / throughput / availability / single-thread)?
\item For \emph{every} block (core, predictor, L1, L2/L3, memory, threads) did I give a choice \emph{and} a reason tied to a constraint?
\item Did I give at least one number (AMAT, CPI, stalls/instr, speedup)?
\item Did I check feature dependencies (Section above)? (OoO $\Rightarrow$ non-blocking; speculation $\Rightarrow$ ROB+predictor; SMT $\Rightarrow$ OoO.)
\item Did I use \textbf{global} miss rate for L2/L3?
\item Did I state costs (area, power, clock, complexity) of each choice?
\item Did I list rejected alternatives and why?
\item Is my issue width within 3--6, unless I justified otherwise?
\item Did I mention power/energy where relevant (dynamic $\tfrac12CV^2f$, DVFS, power gating)?
\item If the scenario hints at compatibility, did I avoid VLIW (or justify)?
\end{enumerate}

\appendix
\chapter{Formula and Number Sheet}\label{app:formulas}
Everything numeric you may need to quote, in one place. $\dagger$ = standard formula not spelled out in the slide text.

\section{Performance, power, cost, dependability}
\begin{center}
\small
\begin{tabular}{L{5.4cm}L{10cm}}
\toprule
\tblhead{\textbf{Item} & \textbf{Formula}}\\
\midrule
CPU time & $\text{IC}\times\CPI\times T_{clk}$ \\
Average CPI & $\sum_i(\text{IC}_i\times\CPI_i)/\text{IC}$ \\
Speedup & $\text{Time}_{old}/\text{Time}_{new}$ \\
Amdahl$^\dagger$ & $1/[(1-f)+f/s]$ \\
Dynamic energy & $\tfrac12 C V^2$ \\
Dynamic power & $\tfrac12 C V^2 f$ \\
Static power & $I_{static}\times V$ (reduce by power gating) \\
Die yield$^\dagger$ & $\text{wafer yield}\times(1+\text{defects/area}\times\text{die area})^{-N}$, $N=11.5$--$15.5$ \\
MTBF; availability & $\text{MTBF}=\text{MTTF}+\text{MTTR}$; $\text{Availability}=\text{MTTF}/\text{MTBF}$ \\
Price vs.\ volume & $-10\%$ per doubling of volume \\
\bottomrule
\end{tabular}
\end{center}

\section{Pipeline and branches}
\begin{center}
\small
\begin{tabular}{L{6.2cm}L{9.2cm}}
\toprule
\tblhead{\textbf{Item} & \textbf{Formula / value}}\\
\midrule
Branch stall cycles per instruction & $f_{branch}\times\text{penalty}$ \\
With prediction$^\dagger$ & $\CPI=\CPI_{base}+f_{branch}\times(1-\text{accuracy})\times\text{penalty}$ \\
5-stage MIPS branch penalty & MEM-resolve 3; EX-resolve 2; ID early-eval 1 (+1 if ALU feeds branch, +2 if load feeds branch) \\
R4000 & 8 stages; load delay 2; branch delay 3 \\
FP unit latency / initiation & ALU 0/1; load 1/1; FP add 3/1; FP mul 6/1; FP div 24/25 \\
Latencies between dependent ops & FP ALU$\to$FP ALU 3; FP ALU$\to$store 2; load$\to$FP ALU 1; load$\to$store 0; int$\to$int 0 \\
1-bit BHT loop (9 taken/1 not) & 80\% accuracy (2 misses per loop execution) \\
2-bit BHT & 1 miss per loop; 4K entries: 82--99\% accuracy \\
$(m,n)$ predictor storage & $2^m$ BHTs of $n$-bit counters; index $=$ PC bits $\Vert$ $m$ history bits \\
$(2,2)$ with 64 entries & 6-bit index = 2 history + 4 PC bits; 16 entries per BHT \\
Delay slots & Compiler fills $\approx$50\% usefully \\
\bottomrule
\end{tabular}
\end{center}

\section{ILP}
\begin{center}
\small
\begin{tabular}{L{6.2cm}L{9.2cm}}
\toprule
\tblhead{\textbf{Item} & \textbf{Value}}\\
\midrule
FP loop (slides) & 9 cycles/iter $\to$ 6 (reordered) $\to$ unroll$\times$4 naive 27 cycles (6.75/iter) $\to$ scheduled 14 cycles (3.5/iter) \\
Scoreboard vs.\ Tomasulo (6-instr.\ example) & Last write: cycle 62 vs.\ 57 \\
Window size & Tomasulo $\le14$; scoreboard $\le5$ \\
Practical issue width & 3--6; window 32--126 \\
Independent instructions needed & Pipeline depth $\times$ number of FUs ($\approx$15--20) \\
Branch frequency (rule) & 1 in 5 instructions \\
VLIW register ports (example) & 7 read/3 write int; 5 read/3 write FP \\
Doubling issue 6$\to$12 needs & 3--4 memory accesses/cycle, 2--3 branches/cycle, rename $>20$ registers/cycle, fetch 12--24 instr./cycle \\
Power5 model & 4 issue/clk; 48 int + 40 FP rename regs; 2--6\% mispredict; 64\,KB I, 32\,KB D, 1.92\,MB L2, 36\,MB L3 \\
\bottomrule
\end{tabular}
\end{center}

\section{Memory}
\begin{center}
\small
\begin{tabular}{L{6.2cm}L{9.2cm}}
\toprule
\tblhead{\textbf{Item} & \textbf{Formula / value}}\\
\midrule
AMAT & $t_{hit}+m\times P$ \\
Two-level AMAT & $t_{L1}+m_{L1}(t_{L2}+m^{loc}_{L2}P_{L2})$; global $m_{L2}=m_{L1}m^{loc}_{L2}$ \\
Stalls/instruction (two-level) & $\text{misses}_{L1}/\text{instr}\times t_{L2}+\text{misses}_{L2}/\text{instr}\times P_{L2}$ \\
Slide example & $1+0.04(10+0.5\times100)=3.4$; stalls $(3.4-1)\times1.5=3.6$ \\
L2 associativity example & DM $10+0.25\times100=35$; 2-way $10+0.20\times100=30$ \\
CPI with memory & $\CPI_{base}+\text{refs/instr}\times m\times P$ (e.g.\ $2+1.2\times0.03\times64=4.3$) \\
Block-size penalty model & 80 cc + 2 cc per 16 B: 82/84/88/96/112 for 16/32/64/128/256 B \\
Best block per cache (AMAT) & 4K: 32 B; 16K, 64K, 256K: 64 B \\
Associativity clock factors & $\times1.10$ (2-way), $\times1.12$ (4-way), $\times1.14$ (8-way) \\
2:1 rule & 2-way of size $N/2\approx$ DM of size $N$ (caches $<128$\,KB) \\
Victim cache & 20--95\% miss reduction, 4\,KB DM, 1--5 entries \\
Way prediction accuracy & $>85\%$ (Alpha 21264, 2-way I-cache) \\
Prefetch AMAT & $t+m\,p\,t_p+m(1-p)P$; example $2.41525$ \\
Stream buffers & ISB 1/4/16 blocks catch 15--25\%/50\%/72\% of I-misses \\
Loop interchange & $1{,}048{,}576\to65{,}536$ (per array); $-93.75\%$ \\
Loop fusion & $262{,}144\to196{,}608$ ($-25\%$) \\
Blocking & $2N^3+N^2\to2N^3/B+N^2$; 3 tiles of $B^2$ words must fit \\
Interleaving & bank = addr mod $\#$banks; banks $\ge$ cycles to access a word \\
Memory-width CPI example & 1-word 4.30; 2-word: simple 5.07 / interleaved 3.63 / wide 3.54; 4-word: 5.69 / 3.09 / 3.84 \\
TLB & 128--256 entries typical; fully associative on high-end \\
Page size / VA split & 4--8\,KB; 32-bit VA = 20-bit VPN + 12-bit offset \\
DRAM refresh & $\sim$8\,ms, $\sim$1\% of time \\
\bottomrule
\end{tabular}
\end{center}

\section{Threads and multicore}
\begin{center}
\small
\begin{tabular}{L{6.2cm}L{9.2cm}}
\toprule
\tblhead{\textbf{Item} & \textbf{Value}}\\
\midrule
Pentium 4 SMT & 1.01 (int rate), 1.07 (FP rate); pairs 0.90--1.58, avg 1.20 \\
Power5 SMT & 1.23 (int rate), 1.16 (FP rate); 0.89--1.41 for two copies \\
Power5 core size & $+24\%$ vs.\ Power4; physical regs 152$\to$240; L2 1.44$\to$1.92\,MB \\
SMT threads per core & 2, 4, or sometimes 8 (Power5: 2) \\
Hardware thread switch & vs.\ 100s--1000s of clocks for a process switch \\
General SMT gain & up to 30\% \\
\bottomrule
\end{tabular}
\end{center}

\chapter{Source Map, Gaps and Glossary}\label{app:src}

\section{Which deck covers what}
\small
\begin{longtable}{L{4.2cm}L{3.0cm}L{8.2cm}}
\toprule
\tblhead{\textbf{Deck (file name)} & \textbf{Slides} & \textbf{Content}}\\
\midrule
\endhead
\texttt{ca01 (1).pdf} & 28 & Course outline; classes of computers; parallelism; Flynn; technology trends; power/energy; cost; dependability; performance equation \\
\texttt{ca03\_noor.pdf} & 88 & Advanced pipelining intro; R4000; FP loop hazards; unrolling; dependences; scoreboard (full worked example); Tomasulo intro; register renaming; VLIW vs.\ superscalar (``ca03'') \\
\texttt{ca03\_noorNew.pdf} & 19 & FP pipeline (latencies, initiation intervals), RISC-V FP, R4000 integer and FP pipeline (``ca03N'') \\
\texttt{ca04\_noorNew.pdf} & 109 & Branch prediction (static, delayed branch, BHT, correlating, two-level, gshare, BTB); hardware speculation, ROB; multiple issue; VLIW; limits of multi-issue and ILP; Superscalar vs.\ VLIW (``ca04'') \\
\texttt{ca05\_noor (1).pdf} & 102 & Cache performance (all techniques), summary table; main memory; DDR; memory hierarchy by device; virtual memory, TLB (``ca05'') \\
\texttt{ca06\_noor.pdf} & 36 & Limits to ILP; multithreading (fine, coarse, SMT); Power5; SMT performance; head-to-head processors (``ca06'') \\
\texttt{ca08\_noorNew.pdf} & 30 & Multicore, SMT, multicore vs.\ SMT, memory hierarchy in multicore (``ca08'') \\
\texttt{lec06-dynasched.pdf} & 69 & Tomasulo vs.\ scoreboard; Tomasulo examples; loop example; speculative Tomasulo with reorder buffer (Kubiatowicz CS252) \\
\texttt{CachePerformance2.pdf} & 54 & Way prediction; pseudo-associative cache; loop interchange/fusion/blocking numeric examples \\
\texttt{Cache-Misses.pdf} & 10 & Taxonomy of cache misses: cold, capacity, conflict, coherence (true/false sharing) \\
\bottomrule
\end{longtable}
\normalsize

\section{Gaps: topics in the course outline but \emph{not} in the provided slides}
If the exam touches these, supplement from the textbook (Hennessy \& Patterson):
\begin{itemize}
\item Vector architectures, SIMD multimedia extensions, GPUs (only named as DLP).
\item Centralised shared-memory vs.\ distributed shared-memory / directory-based coherence, snooping protocols (MSI/MESI), synchronisation, memory consistency models. (Only \emph{coherence misses} are covered.)
\item Warehouse-scale computers, request-level parallelism, cloud computing, physical infrastructure.
\item Detailed treatment of non-blocking cache hit-under-miss statistics beyond the chart; DRAM internals beyond RAS/CAS; flash/disk details.
\end{itemize}

\section{Known inconsistencies in the slides (avoid losing marks)}
\begin{itemize}
\item \textbf{R4000 load delay:} one slide title says ``x1'', the table beneath it shows \emph{two} stall cycles, and the older deck says ``TWO cycle load latency''. Use \textbf{2}.
\item \textbf{Correlating-predictor storage:} the slide compares a 2-bit predictor with 4K entries to a $(2,2)$ predictor with 1K entries and says ``the same number of total bits''; state the claim as ``same hardware budget'' and avoid quoting bit counts.
\item \textbf{Memory-width CPI example:} the slide labels buses ``64-bit'' and ``128-bit'' but the arithmetic treats them as 1-word and 2-word wide; quote the CPI values (4.30, 5.07, 3.63, 3.54, 5.69, 3.09, 3.84) rather than the labels.
\item \textbf{Blocking numbers} ($\approx1.07$ billion $\to\approx2.1$ million) come from a \emph{simplified} model (one miss per column access, tile reuse); present them as illustrative.
\item \textbf{Predict-taken} is described as useless in the base MIPS pipeline only because the target isn't known before the outcome; with a BTB it is meaningful.
\end{itemize}

\section{Glossary}
\small
\begin{longtable}{L{2.4cm}L{13cm}}
\toprule
\tblhead{\textbf{Term} & \textbf{Meaning}}\\
\midrule
\endhead
AMAT & Average memory access time $=t_{hit}+$ miss rate $\times$ miss penalty. \\
BHT/BHR/PHT & Branch history table / branch history register / pattern history table. \\
BTB & Branch target buffer: cache of predicted branch targets (tag + target + optional state). \\
BTFNT & Backward-taken, forward-not-taken static prediction. \\
CDB & Common data bus: broadcast result bus (data + source tag) in Tomasulo. \\
CMP & Chip multiprocessor (multicore). \\
DLP / ILP / TLP / RLP & Data / instruction / thread / request-level parallelism. \\
DVFS & Dynamic voltage and frequency scaling. \\
EPIC & Explicitly parallel instruction computing (IA-64/Itanium, VLIW-style). \\
ISB & Instruction stream buffer (prefetch buffer). \\
MIMD / SIMD / SISD / MISD & Flynn taxonomy classes. \\
OoO & Out-of-order execution. \\
PID & Process identifier used in tags of virtual caches. \\
RAW/WAR/WAW & Read-after-write (true), write-after-read (anti), write-after-write (output) hazards. \\
ROB & Reorder buffer: holds uncommitted results for in-order commit and speculation recovery. \\
RS & Reservation station (Tomasulo buffer with op, operands, tags). \\
SMT & Simultaneous multithreading. \\
TDP & Thermal design power. \\
TLB & Translation lookaside buffer: cache of page-table translations. \\
VLIW & Very long instruction word: compiler packs independent operations into one wide instruction. \\
\bottomrule
\end{longtable}


\end{document}
